% P30 fresh Korean production translation.
% Initial translation authority: ED0005 / NOETH-DE-AUTH-v038-20260805.
% Current scope-rebase authority: ED0011 / NOETH-DE-AUTH-v047-20260808.
% Scope: current German lines 14344--15599.

\providecommand{\mC}{\mathfrak C}
\providecommand{\mE}{\mathfrak E}
\providecommand{\mF}{\mathfrak F}
\providecommand{\mG}{\mathfrak G}
\providecommand{\mH}{\mathfrak H}
\providecommand{\mM}{\mathfrak M}
\providecommand{\mN}{\mathfrak N}
\providecommand{\mO}{\mathfrak O}
\providecommand{\mX}{\mathfrak X}
\providecommand{\mc}{\mathfrak c}
\providecommand{\md}{\mathfrak d}
\providecommand{\me}{\mathfrak e}
\providecommand{\mf}{\mathfrak f}
\providecommand{\mm}{\mathfrak m}
\providecommand{\mn}{\mathfrak n}
\providecommand{\mq}{\mathfrak q}
\providecommand{\mt}{\mathfrak t}
\providecommand{\mv}{\mathfrak v}
\providecommand{\mpideal}{\mathfrak p}

\section*{30. 대수적 수체와 함수체에서 아이디얼론의 추상적 구성.}
\noindent\emph{Math. Ann. 96 (1927), 26--61쪽.}

\setcounter{footnote}{0}

\begin{center}
글쓴이\\[0.5em]
괴팅겐의 에미 뇌터.
\end{center}

이하에서는 아이디얼론이 대수적 수체의 모든 정수적 원소가 이루는 환의 아이디얼론과 일치하는 모든 환, 다시 말해 그 아이디얼들이 소아이디얼들의 거듭제곱곱으로 유일하게 표현되는 모든 환을 추상적으로 특징짓는다. 가장 잘 알려진 예로는 한 부정원의 대수함수체의 모든 정수적 원소가 이루는 환도 이에 속한다. 더 일반적으로 여러 부정원의 경우에도, 크로네커를 따라 최고차원 아이디얼로 제한하여 어떤 몫환인 함수영역으로 넘어가면 마찬가지이다.

이 아이디얼론과 완전히 동치인 공리들은 바탕이 되는 가환환에 대하여 다음과 같다\footnote{아이디얼론의 기본개념 정의는 예컨대 E. Noether, \emph{Idealtheorie in Ringbereichen}, Math. Ann. 83 (1921), 24--66쪽(이하 \emph{아이디얼론}으로 인용)을 보라. 가장 본질적인 개념들은 이 논문에서도, 특히 §§ 1, 4, 5에서 짧게 정식화한다. 완전성을 위하여 이 논문의 직접적 목적에는 필요하지 않은 내용도 얼마간 포함하였다.}.

\begin{description}
\item[I. 약수사슬정리:] 각 아이디얼이 앞선 아이디얼의 진약수인 아이디얼 사슬은 유한 단계에서 멈춘다. 다시 말해, 모든 아이디얼 약수사슬에는 그 뒤의 모든 아이디얼이 서로 같아지는 한 지표가 존재한다.
\item[II. 영아이디얼과 다른 각 아이디얼을 법으로 한 배수사슬정리.] 모두 어떤 고정된 비영아이디얼의 약수인 아이디얼들의 사슬에서, 각 아이디얼이 앞선 아이디얼의 진배수이면 그 사슬은 유한 단계에서 멈춘다.
\item[III.] 곱셈의 단위원이 존재한다.
\item[IV.] 환에는 영인자가 없다.
\item[V. 분수체 안에서의 정수적 닫힘:] 분수체의 원소 가운데 환에 대하여 정수적인 것은 모두 그 환에 속한다.
\end{description}

나는 이 공리들로부터 점차 더 강하게 제한되는 아이디얼론을 전개하여 마침내 찾는 이론에 이르고, 역으로 그 이론에서는 이 공리들이 실제로 모두 성립함을 보인다.

내가 \emph{아이디얼론} § 4에서 보였듯이, 약수사슬정리 I로부터 한 아이디얼을 서로 다른 소아이디얼들에 대응하는 유한개의 으뜸아이디얼의 최소공배수로 나타내는 표현이 나온다. 여기서는 이 증명을 짧게 되풀이한다(§ 6). 아이디얼 몫 개념을 쓰면 증명이 단순해지고, 한 가지 잘못도 바로잡아야 하기 때문이다. 원래 증명은 한 곳에서 가정하지 않은 곱셈 단위원의 존재를 사용한다\footnote{P. Urysohn이 이 잘못을 지적하였다. 내가 고친 증명은 여기 제시하는 E. Artin의 증명보다 조금 번거로웠다. Artin은 이미 더 일찍 스스로 이 틈을 메웠고, 나와 독립적으로 아이디얼 몫을 이용한 단순화도 생각해 냈다. 또한 \glqq{}축약표현\grqq{}을 도입하지 않아도 되는 또 하나의 단순화는 W. Krull의 관련 문제에서 가져왔다. W. Krull, \emph{Algebraische Theorie der Ringe III}, Math. Ann. 92 (1924), 181--213쪽, 정리 I.}.

배수사슬정리를 가정하면(§ 7), 각 비영아이디얼에 따른 잉여류환에서 한 소아이디얼은 모든 원소를 포함하는 단위아이디얼을 제외하고는 진약수를 가질 수 없다. 그러면 단위아이디얼에 대응하는 것들을 제외한 이 아이디얼들의 으뜸성분은 유일하게 결정된다. 이 결과는 조금 덜 날카로운 형태로 이미 Masazo Sono에게서 나타난다\footnote{\emph{On the Reduction of Ideals}, Memoirs of the College of Science, Kyoto University, Series A, 7 (1924), 191--204쪽.}. 그가 잉여류환에 합성열이 존재한다고 둔 가정은 이 환에서의 \glqq{}이중사슬정리\grqq{}와 동치임을 마지막에(§ 10) 짧게 보인다. 여기서 이중사슬정리란 비영아이디얼에 국한하지 않고 약수사슬정리와 배수사슬정리가 모두 성립한다는 뜻이다.

공리 III, 곧 단위원의 존재까지 성립하면 단위아이디얼은 대응 소아이디얼로 나타나지 않는다. 따라서 으뜸성분들은 유일하게 결정되고, 쌍마다 서로소이며, 그 최소공배수는 그 곱과 같다. 공리 I, II, III은 대수적 수체의 모든 유한 오더에 성립한다. 이 경우 데데킨트는 이미 Dirichlet--Dedekind, \emph{Zahlentheorie}, 제3판, 제11보유, § 172에서 완전히 전개된 증명 없이도 한 아이디얼을 쌍마다 서로소인 동종 아이디얼들(으뜸성분들)의 곱으로 유일하게 나타내는 표현을 진술하였다.

가정 IV, 곧 영인자가 없다는 사실로부터 영아이디얼 및 단위아이디얼과 다른 한 아이디얼의 서로 다른 거듭제곱들이 모두 서로 다르고 분수체가 존재한다는 것이 따른다. 끝으로 분수체 안에서 정수적으로 닫혀 있다는 공리 V까지 성립하면 으뜸아이디얼들은 --- IV 때문에 유일하게 결정되는 --- 소아이디얼의 거듭제곱이 되고, 이로써 통상적인 분해정리를 얻는다(§ 8). 이 마지막 증명은 수체의 경우 데데킨트가 이미 제3판 § 172에서 정수적 닫힘으로부터 끌어낸 따름정리에 기초한다. 즉 진분수 원소를 정수적 원소들의 몫으로 나타내되, 분자의 제곱조차 분모로 나누어지지 않게 할 수 있다. 후대의 서술에서는 역시 정수적 닫힘의 결과인 훨씬 덜 투명한 일반화된 가우스 정리나, 그에 대응하면서 조금 더 많은 내용을 말하는 모듈정리들이 이 따름정리를 대신한다. 이 정리들은 모두 적용할 때 기저원소들을 전제하지만 여기서는 아이디얼 자체를 다룬다.

역방향(§ 9)에서는 공리 V 이외의 공리들이 거의 직접 나온다. 공리 V는 진분수 원소를 언제나 그 분자의 어떤 거듭제곱도 분모로 나누어지지 않도록 나타낼 수 있음을 보이면 따른다. 이는 애초 V에서 끌어낸 따름정리를 더 강화한 것이다. 그러나 이 증명은 아이디얼 표현으로부터 통상적인 결론들, 곧 고정된 한 아이디얼을 법으로 한 주아이디얼 표현과 분수아이디얼, 즉 체 안의 아이디얼 몫 이론을 이끌어 낸 뒤에야 가능하다. 소아이디얼 거듭제곱에 의한 표현으로부터 정수적 닫힘이 따른다는 사실도 데데킨트에게 이미 알려져 있었음은 제3판 § 172, 522쪽 아래의 주석이 보여 준다.

공리 I--V에 따르면, 일반적으로는 서로 구별되는 네 분해, 곧 쌍마다 서로소인 최소 아이디얼들로의 분해, 서로 상대소인 아이디얼들로의 분해, 극대 으뜸성분들로의 분해, 기약아이디얼들로의 분해가 서로 일치한다. 역으로 이 일치와 공리 I--IV로부터 정수적 닫힘이 나온다. 영인자를 갖는 환에서는 나머지 공리들이 성립하더라도 --- 이때 V는 몫환에서의 정수적 닫힘으로 정의한다 --- 분해들이 반드시 일치하지는 않는다. 유한 오더의 어떤 아이디얼에 따른 잉여류환이 그 예이다(§ 9).

§ 1에서는 한 환에 대하여 정수적인 원소들의 이론을, 정수적 닫힘에서 나오는 데데킨트의 따름정리까지 개관한다. § 2에서는 계수영역이자 곱셈자영역인 이 환을 갖는 선형형식들의 모듈로 사슬정리들이 어떻게 옮겨지는지를 보인다\footnote{이 형태의 모듈정리는 Prüfer에게서 나왔으며(\emph{Neue Begründung der algebraischen Zahlentheorie}, § 2, Math. Ann. 94 (1924), 198--243쪽), 통상적인 기저정리의 이행보다 더 투명하다.}. 이로부터 § 3에서는 대수적 확대체(제1종)의 정수적 원소들로 넘어갈 때 모든 오더에서 공리 I--IV가 보존되고, 모든 정수적 원소로 이루어진 오더에는 V도 성립한다는 결론을 끌어낸다. 대수적 수체에 관해서는 물론 알려진 이 사실 덕분에 처음에 말한 함수체들도 이 체계 아래 들어온다. 특히 여러 부정원의 정수계수 다항식들에서 유도된 함수영역이 공리들을 만족한다는 간단한 증명으로 크로네커의 아이디얼론도 완전히 포괄된다\footnote{따라서 차이는 판별식정리에서 비로소 나타난다. 한 소수를 법으로 한 함수영역의 잉여류환이 비완전체가 되어 제2종 확대가 생길 수 있기 때문이다.}.

그 뒤에 전개되는 아이디얼론은 분류를 위해서만 쓰인 §§ 2, 3을 전제하지 않는다. §§ 4, 5에서는 공리 I--V와 독립적인 기초를 제시한다. 이로써 이론의 어느 부분이 유한성 가정과 무관한지가 원래의 정초에서보다 더 분명해진다.

\begin{center}
{\bfseries\large § 1.\par}
\vspace{0.35em}
{\bfseries\large 정수적 원소의 이론.\par}
\end{center}

영인자가 없고 곱셈 단위원 \(e\)를 갖는(공리 III, IV) 가환환\footnote{잘 알려져 있듯이, 한 집합에 동치관계와 두 연산인 덧셈과 곱셈이 주어져 통상적인 법칙들을 만족하면 환이 정의된다. 덧셈은 결합적이고 가환적이며 유일하게 역산할 수 있고, 곱셈은 결합적이며 --- 가환환이면 가환적이기도 하고 --- 덧셈에 대하여 분배적이다. 여기서 바탕이 되는 상등의 정의가 집합론적 동일성이 아니라면, 각 부분계에서도 같은 상등의 정의를 유지하기 위하여 그러한 부분계(부분환, 모듈, 아이디얼)는 한 원소와 더불어 그와 같은 모든 원소도 포함해야 한다. 마찬가지로 모든 확대에서는 새 상등의 정의가 옛 정의를 포함한다고 가정해야 한다. \glqq{}유한개의 원소\grqq{}, \glqq{}유일한 대응\grqq{} 등의 모든 개념은 이 상등의 정의에 따라 이해한다. 예컨대 \glqq{}정확히 하나의 원소\grqq{}가 있다는 말은 서로 같은 것들을 제외하면 하나뿐이라는 뜻이다. 같은 원소들을 한 류로 묶고 이 류들을 새 원소로 삼으면 언제나 상등에서 집합론적 동일성으로 넘어갈 수 있다. 상등의 정의에 관한 이 주석들은 R. Hölzer에게서 나왔다.} \(\mT\)를 놓고, \(\mT\) 안에 단위원을 포함하는 고정된 부분환 \(\mR\)을 택한다. 모듈, 오더, 정수적이라는 개념은 줄곧 이 고정된 환 \(\mR\)에 관한 것으로 하며, 뒤에서는 간단히 하려고 표기에서 이를 생략한다\footnote{영인자를 허용하는 대신 \(\mR\)에 약수사슬정리를 가정해도 정수적 원소의 이론은 유지되며, § 2에 따라 보조정리도 증명할 수 있다. 그러나 이 형태의 정수적 원소는 뒤에서 쓰지 않으므로 각주에서만 언급한다. 영인자가 있어도 여기서 다루는 정의들은 모두 유지된다. 잘 알려져 있고 쉽게 알 수 있듯이, 그 대부분은 훨씬 약한 가정만을 필요로 한다. 다만 데데킨트의 정식화 \glqq{}어떤 수가 외피를 가지면 정수적이다\grqq{}는 유지되지 않는다. 영인자가 생기면 두 유한 모듈의 몫이 유한하지 않을 수 있기 때문이다. 그래서 여기서는 오더도 모듈 몫으로가 아니라 직접 정의한다.}.
(\(\mR\)의 원소는 라틴 문자로, 일반적인 \(\mT\)의 원소는 그리스 문자로 나타낸다.)

1. \(\mT\)의 원소들로 이루어진 계 \(\mM\)이 두 원소 \(\alpha,\beta\)와 더불어 그 차도 포함하고, \(\alpha\)와 더불어 모든 \(r\in\mR\)에 대한 \(r\alpha\)도 포함하면, 통상대로 이를 \emph{\(\mR\)-모듈}, 줄여서 \emph{모듈}이라 한다. \(\mM\)의 모든 원소가 \(\mN\)에도 들어 있을 때 \(\mM\)은 \(\mN\)으로 나누어진다,
\[
        \mM\equiv 0(\mN),
\]
고 하고, \(\mM\)을 배수, \(\mN\)을 약수라 한다. 원소가 모두 \(\mR\)에 속하는 \(\mR\)-모듈을 \(\mR\)의 \emph{아이디얼}이라 한다.

분명히 \(\mT\)의 임의의 많은 모듈의 교집합, 곧 최소공배수 \([\ldots,\mM_i,\ldots]\)도 다시 모듈이고, \(\mT\) 자체도 모듈이다. 따라서 \(\mT\)의 원소들의 임의의 계 \(\Sigma\)에서 \emph{생성되는 모듈}은 \(\Sigma\)를 포함하는 모든 모듈의 교집합으로서 유일하게 존재한다. 임의의 모듈계의 합, 곧 최대공약모듈 \((\ldots,\mM_i,\ldots)\)은 그 합집합에서 생성되는 모듈이다. 두 모듈의 곱 \(\mM\mN\)은 \(\mu\)가 \(\mM\)의 모든 원소를, \(\nu\)가 \(\mN\)의 모든 원소를 달릴 때 원소 \(\mu\nu\)들에서 생성되는 모듈로 정의한다. 환원소에 대해 결합법칙과 교환법칙이 성립하므로 모듈의 곱도 이를 만족한다. \(\mT\)의 분배법칙에서 모듈에 대한 분배법칙
\[
        (\ldots,\mM_i,\ldots)\mN=(\ldots,\mM_i\mN,\ldots)
\]
도 따른다. 실제로 \(\mT\)의 바로 그 분배법칙에 따라 양변은 모두 모든 원소 \(\mu_i\nu\)에서 생성되는 모듈이다.

\(\mR\) 자체가 모듈이므로, \(\mR\)이 곱셈자영역이라는 조건은 \(\mR\mM\equiv0(\mM)\)으로 나타낼 수 있다. 가정에 따라 \(\mR\)이 단위원을 포함하므로 \(\mM\equiv0(\mR\mM)\)도 성립하고, 따라서 \(\mM=\mR\mM\)이다.

\(\mM\)이 \emph{유한}하다는 것은 \(\mM\)을 생성하는 원소가 \emph{유한개}인 계 \(\Sigma\)가 적어도 하나 있다는 뜻이다. \(\Sigma\)의 원소 \(\mu_1,\ldots,\mu_r\)를
\[
        \mM=(\mu_1,\ldots,\mu_r)
\]
의 모듈기저라 한다. 두 유한 모듈의 합과 곱, 따라서 유한개의 유한 모듈의 합과 곱도 다시 유한하다. 각각 기저원소들의 합집합과 곱들에서 생성되기 때문이다. 곱에 대한 말은 \(\mM\)의 모든 원소 \(\mu\)가 기저에 따라 \(r_1\mu_1+\cdots+r_k\mu_k\)로 표현된다는 사실\footnote{\(\mR\)이 단위원을 포함하지 않으면 기저표현은 \(r_1\mu_1+\cdots+r_k\mu_k+n_1\mu_1+\cdots+n_k\mu_k\) 꼴이 된다. 여기서 정수 \(n\)은 환원소가 아니라 반복된 덧셈 또는 뺄셈을 나타내는 기호이다.}과 \(\mT\)의 곱셈 교환법칙에서 따른다. 이 두 사실은 여기서 처음으로 쓰인다.

2. \(\mR\)의 원소를 모두 포함하면서 \(\mT\) 안에 놓인 \(\mR\)의 확대환을 \emph{\(\mR\)-오더}, 줄여서 \emph{오더}라 한다. \(\mT\) 안의 임의의 많은 오더의 교집합도 다시 오더이고 \(\mT\)도 오더이므로, \(\mT\)의 원소들의 임의의 계 \(\Sigma\)에서 생성되는 오더가 유일하게 존재한다. 오더의 합과 곱은 모듈의 경우와 마찬가지로 정의할 수 있고, 특히 모든 오더에 대해 \(\mO^2=\mO\)이다.

모든 오더는 동시에 \(\mR\)-모듈이다. 오더가 유한하다는 것은 그것이 유한 모듈이라는 뜻이다. 합과 곱은 모듈의 합과 곱으로 만들어지므로 1에 따라 유한 오더의 합과 곱도 다시 유한 오더이다. 또한 다음 \emph{추이법칙}이 성립한다. \(\mA\)가 유한 \(\mR\)-오더이고 \(\mB\)가 유한 \(\mA\)-오더이면, \(\mB\)는 유한 \(\mR\)-오더이기도 하다. 실제로 \(\mB\)는 동시에 \(\mR\)-오더, 따라서 \(\mR\)-모듈이 된다. 그런데 \(\alpha_1,\ldots,\alpha_\rho\)가 \(\mR\)에 관한 \(\mA\)의 모듈기저이고 \(\beta_1,\ldots,\beta_\sigma\)가 \(\mA\)에 관한 \(\mB\)의 모듈기저이면, 곱 \(\alpha_i\beta_k\)들이 분명 \(\mR\)에 관한 \(\mB\)의 모듈기저를 이룬다.

3. \(\mT\)의 원소 \(\alpha\)에서 생성되는 오더 \(\mR_\alpha\)가 \emph{유한}하면 \(\alpha\)를 \emph{\(\mR\)에 대하여 정수적}, 줄여서 \emph{정수적}이라 한다. 달리 말하면 모듈의 약수사슬
\[
        \mA_\mu=(\alpha^0,\alpha,\ldots,\alpha^{\mu-1})
\]
이 유한 단계에서 멈추어 \(\mA_n=\mA_{n+1}=\cdots=\mA_{n+\nu}=\cdots\)가 된다는 뜻이다.

이 정의는 4에서 쓸 다음 두 번째 형태로도 쓸 수 있다. \(\mT\)의 원소 \(\alpha\)가 정수적이라는 것은 영모듈과 다른 주모듈 \(\mC\), 곧 한 원소 \(\gamma\ne0\)에서 생성되는 \(\mC\)가 적어도 하나 존재하여 오더 \(\mR_\alpha\)와 \(\mC\)의 곱이 유한 모듈이 된다는 뜻이다. 달리 말하면 모듈의 사슬 \(\mC\mA_\mu\)가 유한 단계에서 멈춘다\footnote{\(\mR\)에 영인자를 허용하면 정칙 주모듈의 존재를 요구해야 한다. 즉 \(\gamma\)가 영인자가 아닌 정칙원소라고 가정해야 하며, 영인자가 없는 환에서는 이는 \(\gamma\ne0\)과 같다.}.

끝으로 이 정의는 통상적인 다음 형태와도 같다. \(\mT\)의 원소 \(\alpha\)가 정수적이라는 것은 \(\mR\)의 원소 \(r_i\)를 계수로 하는 방정식
\[
        \alpha^n+r_1\alpha^{n-1}+\cdots+r_n=0
\]
을 적어도 하나 만족한다는 뜻이다.

이 여러 정의는 실제로 서로 같다. \(\mR_\alpha\)는 \(\alpha\)의 모든 거듭제곱에서 생성되는 모듈과 같으므로, 유한한 \(\mR_\alpha\)에는 언제나 이 거듭제곱들로 이루어진 모듈기저를 택할 수 있다. 예컨대 그 기저가 \(e=\alpha^0,\alpha,\ldots,\alpha^{n-1}\)이면 \(\mA_\mu\)의 사슬이 \(\mA_n\)에서 멈춘다는 뜻이고, 역으로 \(\mA_n=\mA_{n+\nu}=\cdots\)에서 이 기저가 나온다. \(\mA_\mu\)의 사슬이 멈춘다는 것은 방정식
\[
        \alpha^n+r_1\alpha^{n-1}+\cdots+r_n=0
\]
이 성립하는 것과도 같다. \(\mR_\alpha\)가 유한하면 \(\mC\mR_\alpha\)도 유한하므로 \(\mC\mA_\mu\)의 사슬이 멈춘다. 역으로 \(\mC\mR_\alpha\)가 유한하여 \(\mC\mA_\mu\)가 멈추면 \(\mA_\mu\)도 멈추고, 따라서 \(\mR_\alpha\)가 유한하다. \(\mC\)를 비영 주모듈로 가정했으므로 \(\mC\mA_n=\mC\mA_{n+\nu}\)에서 언제나 \(\mA_n=\mA_{n+\nu}\)가 나오기 때문이다.

정수적 원소들이 환을 이룬다는 사실과 정수적 종속성의 추이법칙은 다음 보조정리에 기초한다.

\medskip
\noindent\textbf{보조정리.} \(\mS\)가 \(\mT\) 안의 유한 오더이고 \(\alpha\)가 \(\mS\)의 임의의 원소이면, \(\alpha\)에서 생성되는 오더 \(\mR_\alpha\)는 유한하다. 다시 말해, 유한 오더의 모든 원소는 정수적이다.

증명은 통상적인 추론에서 나온다. \(\sigma_1,\ldots,\sigma_n\)이 \(\mS\)의 모듈기저이면, 환의 성질에 따라 \(\alpha\sigma_i\)도 \(\mS\)에 속한다. 따라서
\[
        \alpha\sigma_i=r_{i1}\sigma_1+\cdots+r_{in}\sigma_n,
\]
이고, 여기서
\[
        \left|r_{ik}-\alpha e_{ik}\right|=0,
\]
을 얻는다. \(i\ne k\)이면 \(e_{ik}=0\), \(e_{ii}=e\)이므로 이 식은 \(\alpha\)가 정수적임을 보인다. 행렬식이 사라짐을 보일 때 \(\mT\), 따라서 \(\mS\)에 영인자가 없다는 가정을 쓴다\footnote{\(\mR\)에 약수사슬정리를 가정하되 영인자를 허용하면, 이 보조정리는 § 2의 모듈정리에서 직접 따른다. 그 정리에 따르면 사슬정리가 \(\mS\)의 모듈들로 옮겨진다. 수체에 대한 이 증명 역시 데데킨트에게서 나왔고(제4판, § 173, III), 문헌에서 사슬정리를 처음 적용한 예이다. 아래 4의 따름정리들에서 데데킨트는 그 대신 대수적 수체의 한 아이디얼에 따른 잉여류의 수가 유한하다는 더 특수한 사실을 아직 사용한다.}.

\(\mT\) 안에서 정수적인 모든 원소의 계 \(\mS\)는 오더를 이룬다. \(\mS\)는 \(\mR\)을 포함한다. 실제로 \(\mR\)의 원소들은 정수적이다. \(\mR\)은 단위원을 기저로 하는 유한 \(\mR\)-모듈이기 때문이다. \(\mS\)는 환의 성질도 갖는다. 임의의 두 정수적 원소 \(\alpha,\beta\)에서 생성되는 오더 \(\mR_{\alpha,\beta}\)는 곱 \(\mR_\alpha\cdot\mR_\beta\)와 같아서 유한하다. 따라서 보조정리에 의해 유한 오더의 원소 \(\alpha\pm\beta\)와 \(\alpha\beta\)는 정수적이다.

\medskip
\noindent\textbf{정수적 종속성의 추이법칙.} \(\mS\)가 정수적 원소들로 이루어진 오더이면, \(\mS\)에 대하여 정수적인 \(\mT\)의 모든 원소는 \(\mR\)에 대해서도 정수적이다. 정의에 따라 \(\alpha\)는 방정식
\[
        \alpha^m+\sigma_1\alpha^{m-1}+\cdots+\sigma_m=0
\]
을 만족한다. 그러므로 \(\alpha\)는 \(\sigma_1,\ldots,\sigma_m\)에서 생성되는 오더 \(\overline{\mS}=\mR_{\sigma_1,\ldots,\sigma_m}\)에 대해서도 정수적이다. 이 오더는 곱 \(\mR_{\sigma_1}\cdot\mR_{\sigma_2}\cdots\mR_{\sigma_m}\)이므로 유한하다. 따라서 2의 추이법칙에 따라 \(\alpha\)에서 생성되는 유한 \(\overline{\mS}\)-오더 \(\overline{\mS}_\alpha\)는 유한 \(\mR\)-오더이기도 하고, 보조정리에 따라 \(\alpha\)는 유한 오더의 원소로서 정수적이다.

4. \(\mR\)에 대하여 정수적인 \(\mT\)의 모든 원소가 \(\mR\)에 속하면 \(\mR\)을 \emph{\(\mT\) 안에서 정수적으로 닫혀 있다}고 한다. 3의 정수적 원소 정의의 두 번째 형태에 따르면, \(\mT\)의 원소 \(\alpha\)가 \(\mR\)에 속할 필요충분조건은 영모듈과 다른 주모듈 \(\mC\)가 적어도 하나 존재하여 \(\mC\cdot\mR_\alpha\)가 유한 모듈을 이루는 것이다\footnote{\(\mR\)에 영인자를 허용하면 \(\mC\)를 다시 정칙 주모듈로 가정해야 하고, 따름정리 I의 원소 \(c\)도 정칙이라고 가정해야 한다.}.

데데킨트는 정수적으로 닫힌 환이라는 개념에서 아이디얼론에 이 개념을 적용할 수 있게 하는 중요한 따름정리 두 개를 끌어냈다(제3판, § 172).

\medskip
\noindent\textbf{데데킨트의 따름정리 I.} \emph{\(\mR\)이 \(\mT\) 안에서 정수적으로 닫혀 있고, 추가로 \(\mR\)의 아이디얼에 대한 약수사슬정리(공리 I)가 성립한다고 하자. \(\beta\)가 \(\mT\)의 정수적이지 않은 원소이고 \(c\ne0\)이 \(\mR\)의 원소이면, 원소들의 열
\[
        c,\;c\beta,\ldots,c\beta^r,\ldots
\]
에서 \(c,\ldots,c\beta^\sigma\)는 정수적이고 그 뒤의 모든 원소는 정수적이지 않게 하는 지수 \(\sigma\ge0\)가 존재한다.}

모든 원소 \(c\beta^r\)가 정수적이라면 \(\mR\)의 정수적 닫힘에 따라 모두 \(\mR\)에 속한다. 그러면 모듈의 사슬 \(\mC\mB_1,\mC\mB_2,\ldots,\mC\mB_r,\ldots\), 여기서 \(\mC\)는 \(c\)에서, \(\mB_r\)은 \(\beta^0,\ldots,\beta^{r-1}\)에서 생성되는 모듈인데, 이는 \(\mc_r=(c,c\beta,\ldots,c\beta^{r-1})\)인 아이디얼의 사슬 \(\mc_1,\ldots,\mc_r,\ldots\)로 바뀐다. 이 사슬은 가정에 따라 유한 단계에서 멈춘다. 그러면 \(\mR_\beta\)가 유한하고 \(\beta\)가 정수적이 되어 가정에 모순된다. 따라서 \(c\beta^r\) 가운데 정수적이지 않은 원소가 있다. 더 나아가 하나가 정수적이지 않으면 뒤의 모든 원소도 정수적이지 않다. 실제로 \(c\beta^\tau\)가 정수적, 따라서 \(\mR\)에 속한다면 \(\rho<\tau\)인 \(c\beta^\rho\)는 방정식
\[
        (c\beta^\rho)^\tau-c^{\,\tau-\rho}(c\beta^\tau)^\rho=0
\]
을 만족하므로 역시 정수적이다. 따라서 \(c\beta^{\sigma+1}\)이 최초의 정수적이지 않은 원소라면 \(c\)가 정수적이므로 \(\sigma\ge0\)이고, 이것이 찾는 지수이다.

확대환 \(\mT\)를 \(\mR\)의 분수체, 곧 모든 원소쌍을 첨가하여 생기는 체\footnote{잘 알려진 Steinitz의 정식화, J. f. M. 137. \(\mR\)에 영인자를 허용하면, 분수체 대신 \(\mR\)의 정칙원소를 분모로 하는 모든 몫쌍을 첨가하여 얻는 몫환을 써야 한다. 이 경우에는 데데킨트의 따름정리 II가 더 이상 성립하지 않는다. 일반적으로 \(n\)이 정칙원소가 되지 않기 때문이다.}로 특수화하면 다음이 따른다.

\medskip
\noindent\textbf{데데킨트의 따름정리 II.} \emph{\(\mR\)이 그 분수체 \(\mK\) 안에서 정수적으로 닫혀 있고 \(\mR\)에서 아이디얼에 대한 약수사슬정리가 성립한다고 하자. \(\beta\)가 \(\mK\)의 정수적이지 않은 원소이면, \(\beta\)를 \(\mR\)의 원소들의 몫
\[
        \beta=m/n
\]
으로 나타내되 \(m^2/n\)도 정수적이지 않게 할 수 있다.}

\(\beta=b/c\), \(c\ne0\)이라 쓰면 열 \(c,c\beta=b,c\beta^2,\ldots\)의 첫 두 원소는 확실히 정수적이다. 따라서 따름정리 I의 지수 \(\sigma\)는 \(\sigma\ge1\)이다. \(m=c\beta^\sigma\), \(n=c\beta^{\sigma-1}\)로 놓으면 --- \(\sigma\ge1\)이므로 이들은 열의 정수적 원소들이다 --- \(m/n=\beta\)이고 \(m^2/n=c\beta^{\sigma+1}\)은 정수적이지 않다.

\emph{정수적 닫힘의 추이법칙}은 다음 형태로 성립한다. \(\mR\)을 \(\mT\)의 임의의 환이라 하고, \(\mS\)를 \(\mT\) 안에서 \(\mR\)에 대하여 정수적인 모든 원소의 계라 하자. 그러면 \(\mS\)는 \(\mS\)에 대하여 정수적인 \(\mT\)의 모든 원소로도 이루어진다. 실제로 \(\mS\)에 대하여 정수적인 모든 원소는 3의 추이법칙에 따라 \(\mR\)에 대해서도 정수적이므로 \(\mS\)에 속한다.

\subsection*{§ 2. 유한 모듈영역에서의 사슬정리}

사슬정리들이 옮겨진다는 모듈정리는 여기서의 적용에서는 확대환의 모듈에 대해서만 나타난다. 그러나 일반 모듈영역의 경우에도 증명이 같으므로 그러한 일반 영역을 바탕에 둔다.

원소 \(\alpha,\beta,\ldots\)의 계 \(M\)에 다음 두 연산, 곧 원소들의 덧셈 연산과 원소에 \(\mR\)의 원소를 곱하는 연산이 주어져 언제나 \(M\)의 원소를 유일하게 낳고, \(M\)이 덧셈에 관하여 아벨군을 이루며 곱셈에는 결합법칙이 성립하고 분배법칙이 두 형태로 모두 성립하면, \(M\)을 \emph{환 \(\mR\)에 관한 모듈영역}이라 한다\footnote{\emph{아이디얼론}, § 9를 보라.}.

모듈영역에는 § 1, 1에서 제시한 모듈 정의 가운데 곱셈을 언급하지 않는 것들이 분명 모두 유지된다. 특히 \(M\)에서 유한개의 원소 \(\xi_1,\ldots,\xi_k\)를 택하여 \(M\)이 이 원소들에서 생성되는 모듈, 곧 \(\mR\)에 단위원이 있을 때 모든 선형형식 \(r_1\xi_1+\cdots+r_k\xi_k\)의 계와 같으면 \(M\)을 \emph{유한 모듈영역}이라 한다. 단위원이 없으면 여기에 항 \(n_i\xi_i\)들도 더해진다(각주 8 참조).

\medskip
\noindent\textbf{모듈정리.} \emph{\(M\)이 단위원을 갖는 가환환 \(\mR\)에 관한 유한 모듈영역이고, \(\mR\)에서 아이디얼의 약수사슬정리 또는 배수사슬정리가 성립하면, \(M\)에서 모듈의 약수사슬정리 또는 배수사슬정리가 각각 성립한다.}

증명은 \(M\)의 각 모듈에 \(\mR\)의 유한개의 아이디얼로 이루어진 한 계를 유일하게 대응시키되, 모듈의 진약수 또는 진배수에는 대응 아이디얼 가운데 적어도 하나의 진약수 또는 진배수가 각각 대응하게 하는 데 기초한다.

\(M\)의 한 원소가 \(\xi_1,\ldots,\xi_i\)의 선형형식으로 적어도 한 가지 방식으로 표현되지만 \(\xi_1,\ldots,\xi_{i-1}\)의 선형형식으로는 표현되지 않으면 그 원소의 \emph{길이가 \(i\)}라고 한다. 이제 \(M\)의 임의의 모듈 \(\mA\)에 \(\mR\)의 \(k\)개 아이디얼 \(\mathfrak a_1,\ldots,\mathfrak a_k\)를 다음과 같이 대응시킨다. \(\mA_i\)는 길이가 \(\le i\)인 \(\mA\)의 모든 원소가 이루는 모듈이고, \(\mathfrak a_i\)는 \(\mA_i\)에서 \(\xi_i\)의 계수들이 이루는 계이다. 우선 다음이 따른다.

\(\mB\)가 \(\mA\)의 \emph{진약수}이고 \(\mathfrak b_1,\ldots,\mathfrak b_k\)가 \(\mB\)에 대응하는 아이디얼들이면, \(\mathfrak b\) 가운데 적어도 하나는 대응하는 \(\mathfrak a\)의 \emph{진약수}이다. 실제로 \(\mA\equiv0(\mB)\)에서 \(\mA_\lambda\equiv0(\mB_\lambda)\), 따라서 \(\lambda=1,2,\ldots,k\)에 대해 \(\mathfrak a_\lambda\equiv0(\mathfrak b_\lambda)\)가 나온다. 가정에 따라 \(\mB\)에는 \(\mA\)에 들어 있지 않은 원소가 있고, 따라서 그런 원소 가운데 최소 길이 \(i\)인 것도 있다. 그러면 \(\mathfrak b_i\)는 \(\mathfrak a_i\)의 진약수이다. 실제로 그러한 최소 길이의 원소
\[
        \beta=b_1\xi_1+\cdots+b_i\xi_i
\]
를 \(\mB\)에서 택하면 \(b_i\equiv0(\mathfrak b_i)\)이지만 \(b_i\not\equiv0(\mathfrak a_i)\)이다. 만일 \(b_i\equiv0(\mathfrak a_i)\)라면 원소
\[
        \alpha=a_1\xi_1+\cdots+b_i\xi_i\equiv0(\mA)
\]
가 존재하고, 따라서 \(\beta\)보다 길이가 짧으면서 \(\beta-\alpha\equiv0(\mB)\), \(\beta-\alpha\not\equiv0(\mA)\)인 원소가 존재하여 최소성에 모순된다.

이제 약수사슬 또는 배수사슬
\[
        \mA_1,\mA_2,\ldots,\mA_r,\ldots
\]
이 주어지고 \(\mR\)에서 각각 약수사슬정리 또는 배수사슬정리가 성립한다고 하자. \(\mathfrak a_{r1},\ldots,\mathfrak a_{rk}\)가 \(\mA_r\)에 대응하는 아이디얼이면, \(\lambda=1,\ldots,k\)마다 열
\[
        \mathfrak a_{1\lambda},\mathfrak a_{2\lambda},\ldots,\mathfrak a_{r\lambda},\ldots
\]
은 약수사슬 또는 배수사슬을 이룬다. 따라서 어떤 지표 \(\mu\)가 존재하여 모든 \(\lambda\)에 대해 아이디얼 \(\mathfrak a_{\mu\lambda}\)가 그 뒤의 모든 아이디얼과 같다. 그러면 위에서 증명한 바에 따라 모듈 \(\mA_\mu\)도 그 뒤의 모든 모듈과 같아진다. 이로써 \emph{사슬정리들이 옮겨졌다}.

곧바로 다음을 얻는다.

\medskip
\noindent\textbf{모듈정리의 따름정리.} \emph{\(\mR\)에서 영아이디얼과 다른 각 아이디얼을 법으로 한 배수사슬정리가 성립하면, \(M\)에서는 대응 아이디얼 \(\mathfrak c_1,\ldots,\mathfrak c_k\)가 모두 비영아이디얼인 각 모듈 \(\mC\)를 법으로 한 배수사슬정리가 성립한다.}
실제로 이 가정 아래에서는 배수사슬 \(\mathfrak a_{1\lambda},\ldots,\mathfrak a_{r\lambda},\ldots\)가 유한 단계에서 멈춘다.

\medskip
\noindent\textbf{덧붙이는 주석.} \(\mR\)이 단위원 없는 환이어도 약수사슬정리와 비영아이디얼을 법으로 한 배수사슬정리는 여전히 옮겨진다. 실제로 \(\mA\) 대신
\[
        a_1\xi_1+\cdots+a_k\xi_k+n_1\xi_{k+1}+\cdots+n_k\xi_{2k}
\]
꼴의 모든 원소가 이루는 계 \(\overline{\mA}\)를 고찰한다. \(\xi_{k+\lambda}\)를 \(\xi_\lambda\)로 바꾸면 이 계는 다시 \(\mA\)가 된다. 위와 마찬가지로 \(\overline{\mA}\)에 \(2k\)개 아이디얼 \(\mathfrak a_1,\ldots,\mathfrak a_k,\mathfrak n_1,\ldots,\mathfrak n_k\)를 대응시킬 수 있다. 여기서 \(\mathfrak a_\lambda\)는 \(\mR\)의 아이디얼이고 \(\mathfrak n_\lambda\)는 정수환의 아이디얼이다. \(\mathfrak n_\lambda\)에는 약수사슬정리와 각 비영아이디얼을 법으로 한 배수사슬정리가 성립하므로, \(\overline{\mA}\), 따라서 \(\mA\)에 대한 약수사슬정리의 증명은 그대로 유지된다. 배수사슬정리에서는 \(\mC\), 또는 \(\overline{\mC}\)에 대응하는 \(2k\)개 아이디얼이 모두 비영아이디얼이라고 요구해야 한다.

\providecommand{\mK}{\mathfrak{K}}
\providecommand{\mL}{\mathfrak{L}}
\providecommand{\mT}{\mathfrak{T}}
\providecommand{\mC}{\mathfrak{C}}
\providecommand{\mF}{\mathfrak{F}}
\providecommand{\ma}{\mathfrak{a}}
\providecommand{\mb}{\mathfrak{b}}
\providecommand{\mc}{\mathfrak{c}}
\providecommand{\md}{\mathfrak{d}}
\providecommand{\me}{\mathfrak{e}}
\providecommand{\mf}{\mathfrak{f}}
\providecommand{\mm}{\mathfrak{m}}
\providecommand{\mn}{\mathfrak{n}}
\providecommand{\mo}{\mathfrak{o}}
\providecommand{\mq}{\mathfrak{q}}
\providecommand{\mt}{\mathfrak{t}}
\providecommand{\mv}{\mathfrak{v}}
\providecommand{\mpideal}{\mathfrak{p}}

\subsection*{§ 3. 유한 확대체로의 이행}

수체와 함수체를 이 체계 아래 넣으려면, 서론에서 정식화한 공리 I--V가 유한 확대체로 넘어갈 때 어떻게 옮겨지는지를 보여야 한다.

1. \emph{\(\mR\)에서 공리 II, 곧 배수사슬정리를 제외한 공리 I--V가 모두 성립한다고 하자. \(\mK\)는 \(\mR\)의 분수체이고 \(\mL\)은 \(\mK\)의 유한 제1종 확대체라 하자\footnote{Steinitz에 따르면, 모든 원소가 어떤 적당한 확대체에서 서로 다른 일차인자들로 분해되는 소함수의 영점이면 그 대수적 확대체를 \glqq{}제1종\grqq{}이라 한다.}. 그러면 \(\mL\) 안에서 \(\mR\)에 대하여 정수적인 모든 원소의 계 \(\mS\)에는 같은 공리들이 성립하며, \(\mS\) 안의 모든 오더에는 정수적 닫힘을 제외한 이 공리들이 모두 성립한다.}

§ 1, 3에 따라 \(\mS\)는 공리 III, IV, 곧 단위원의 존재와 영인자가 없다는 성질이 성립하는 환을 이룬다. § 1, 4의 끝에 따라 \(\mS\)는 \(\mL\) 안에서 정수적으로 닫혀 있다. 또한 \(\mL\)의 모든 원소는 \(\mR\)의 적당한 원소를 곱하면 \(\mS\)의 원소가 되므로 \(\mL\)은 \(\mS\)의 분수체이다. 따라서 공리 V가 성립한다.

I의 증명은 알려진 추론으로 얻어진다. 제1종 확대인 \(\mL\)은 \(\mK\)에 한 원소 \(\alpha\)를 첨가하여 생기며, 위에 따라 \(\alpha\)가 \(\mR\)에 대하여 정수적이라고 가정할 수 있다. \(\alpha\)와 더불어 \(\mK\)에 대한 \(\mL\)의 갈루아체 안에 있는 켤레원소 \(\alpha',\alpha'',\ldots\)도 정수적이며, 따라서 그 차들의 곱과 이 차곱의 제곱 \(D\)도 정수적이다. \(\alpha,\alpha',\ldots\)가 모두 서로 다르므로 \(D\)는 \(\mK\)의 비영원소이고, \(\mR\)이 \(\mK\) 안에서 정수적으로 닫혀 있으므로 \(\mR\)의 원소이다. 통상적인 추론에 따라 \(\mS\)는 \(\mR\)에 관한 유한 모듈영역 \(M\), 곧 원소 \(\xi_i=\alpha^i/D\)들이 생성하는 영역에 포함된다. 이를 보이려면 \(\mS\)의 한 원소를 \(\xi_i\)들로 나타낸 표현에서 켤레원소들로 넘어간 다음 표현계수에 대한 선형방정식계를 풀기만 하면 된다. § 2의 모듈정리에 따라 \(M\) 안의 모든 \(\mR\)-모듈, 특히 \(\mS\)의 모든 아이디얼에 약수사슬정리가 성립한다. \(\mS\) 안의 임의의 오더의 아이디얼들도 \(M\) 안의 \(\mR\)-모듈이므로 이들에도 약수사슬정리가 성립하고, 모든 오더에는 공리 III, IV도 성립한다.

2. \emph{\(\mR\)에서 공리 I--V가 모두 성립하면 \(\mS\)에서도 모두 성립한다. \(\mS\) 안의 모든 오더에는 정수적 닫힘을 제외한 모든 공리가 여전히 성립한다.}

1을 고려하면 공리 II만 더 증명하면 된다. § 2의 모듈정리의 따름정리에 따라 다음을 보이면 충분하다. \(\mS\)의 비영아이디얼을 \(M\) 안의 \(\mR\)-모듈 \(\mC\)로 볼 때, \(\mC\)에 대응하는 \(\mR\)의 아이디얼 \(\mc_i\)가 모두 비영아이디얼이어야 한다. 그런데 \(\mC\)는 \(\mK\)에 관하여 \(M\)과 같은 유한 선형 랭크를 갖는다. 실제로 \(\mC\)는 임의의 한 원소 \(\gamma\)와 더불어 \(\gamma\alpha^i\)도 포함한다. 따라서 각 \(\xi_i\)에 대해 \(c_i\xi_i\)가 \(\mC\)에 속하게 하는 비영원소 \(c_i\in\mR\)가 존재한다. \(\xi_i\)에 의한 표현의 유일성에 따라 --- 지표는 체의 차수보다 작은 값만 달린다 --- \(c_i\)는 \(\mc_i\)에 속하므로 \(\mc_i\)는 비영아이디얼이다. 이 논의는 \(M\)과 같은 랭크를 갖는 모든 오더에도 유지된다. 더 낮은 랭크의 오더에 대해서는 그 분수체로 넘어간다. 이 체는 \(\mK\)와 \(\mL\)의 중간체로서 역시 제1종이고, \(\mK\)에 대한 그 차수는 오더의 선형 랭크와 같으므로 위 논의가 유지된다. 끝으로 \(\mS\)와 다른 \(\mS\)의 오더는 결코 정수적으로 닫혀 있지 않다. 모든 오더가 \(\mR\)도 포함하므로, 정수적으로 닫혀 있다면 \(\mS\)까지 포함해야 하기 때문이다.

그러므로 수체와 함수체를 이 체계 아래 넣으려면 기초영역, 곧 정수환, 한 부정원의 다항식환, 여러 부정원의 다항식환의 함수영역에서 공리 I--V가 성립함을 보이면 충분하다.

3. \emph{유리정수환 \(\mR\), 또는 한 체를 계수로 하는 한 부정원의 다항식환 \(\mR\)에는 공리 I--V가 성립한다.} I과 II는 각 비영수 또는 한 부정원의 각 비영다항식을 법으로 한 잉여류가 각각 유한개, 또는 선형독립인 것이 유한개라는 사실에서 나온다. 또는 모든 아이디얼이 주아이디얼이라는 사실에서도 나온다. 이 사실로부터 아이디얼이 소아이디얼의 거듭제곱곱으로 유일하게 분해되고, 따라서 한 비영아이디얼을 나누는 아이디얼은 유한개뿐이라는 것이 따른다. 공리 III, IV는 분명 성립하고, IV는 상등의 정의에서 따른다. 정수적 닫힘 V는 다음에서 나온다. \(\mR\)의 원소들은 단위원, 곧 1의 약수들을 제외하면 기약원소들로 유일하게 분해된다. 따라서 \(\mR\)에 들어 있지 않은 분수체의 모든 원소를 \(\mR\)의 서로소인 두 원소의 몫으로 나타낼 수 있고, 그 분자의 어떤 거듭제곱도 분모로 나누어지지 않는다. 그러므로 그러한 원소는 \(\mR\)에 대하여 정수적임을 특징짓는 방정식을 만족할 수 없다.

4. \emph{이제 \(\mR\)을 한 체의 계수 또는 유리정수 계수를 갖는 여러 부정원의 모든 다항식이 이루는 환이라 하자.} 이 경우에도 원소가 단위원을 제외하고 기약원소들로 유일하게 분해되므로, 3에서와 같이 공리 III, IV, V가 성립한다. 약수사슬정리 I은 아이디얼기저의 존재에 관한 힐베르트 정리의 직접적인 결과이지만, 배수사슬정리 II는 여기서 성립하지 않는다.

그러나 최고차원 아이디얼로 제한하는 것과 같은 \emph{함수영역}으로 넘어가면 공리 II까지 성립하게 할 수 있다. 그러면 I도 힐베르트 정리를 쓰지 않고 3에서와 같이 증명할 수 있다. 실제로 \(\mR\)에 한 부정원 \(u\)를 첨가하여, \(\mR\)의 원소를 계수로 하는 \(u\)의 모든 다항식이 이루는 환 \(\mR^*\)을 고찰하고 상등을 계수의 상등으로 정의한다. 통상대로 \(\mR^*\)의 한 다항식의 계수들의 최대공약원 --- 여기서는 아이디얼 약수가 아니라 \(\mR\)의 다항식이라는 뜻의 약수이다 --- 이 \(\mR\)의 단위원이면 그 다항식을 \(\mR\)에 관하여 원시적이라 한다. 그러면 \(\mR^*\)의 모든 다항식은 \(\mR\)의 한 원소와 \(\mR^*\)의 한 원시다항식의 곱이고, \(\mR^*\)의 원시다항식들의 곱은 다시 원시적이다.

\emph{\(\mR^*\)의 분수체 원소 가운데 분모가 \(\mR^*\)의 원시다항식인 것들의 전체는 하나의 환, 곧 \(\mR\)의 함수영역 \(\mF\)를 이룬다.} 함수영역 \(\mF\)에서는 분자와 분모가 모두 \(\mR^*\)의 원시다항식인 원소들이 정확히 단위원이다. 따라서 \(\mF\)의 모든 원소는 \(\mR\)의 한 원소와 \(\mF\)의 한 단위원의 곱으로 유일하게 표현된다. 이로써 \(\mR\) 원소들의 분해정리가 \(\mF\)의 원소들로 옮겨지고, \(\mF\)에는 공리 III, IV와 더불어 3에서처럼 V도 성립한다.

또한 \(\mF\)의 모든 아이디얼은 주아이디얼이다. 실제로 한 아이디얼은 \(\mF\)의 임의의 한 원소와 더불어, 여기에 한 단위원을 곱하여 얻는 \(\mR\)의 원소도 포함한다. 그리고 \(\mR\)의 임의의 두 원소와 더불어 그 최대공약원도 포함한다. 왜냐하면 아이디얼이 \(f(x)=t(x)\overline f(x)\)와 \(g(x)=t(x)\overline g(x)\)를 포함하면 \(t(x)(\overline f(x)+u\overline g(x))\)도 포함하고, \(\overline f,\overline g\)가 서로소여서 그 선형결합이 단위원이 되면 \(t(x)\)도 포함하기 때문이다. 따라서 아이디얼의 임의의 원소 \(f(x)\)에서 시작하여, 아이디얼 안에 \(f(x)\)로 나누어지지 않는 \(g(x)\)가 있으면 최대공약원 \(t(x)\)도 아이디얼에 속하며 \(f(x)\)의 진약수가 된다. 이 과정을 유한번 되풀이하면 아이디얼의 한 기저원소에 이르고, 따라서 그 아이디얼은 주아이디얼이다. 그러므로 \(\mF\)에서 아이디얼은 \(\mR\) 원소들의 분해에 대응하여 소아이디얼들의 거듭제곱곱으로 유일하게 분해되고, 이로부터 3에서와 같이 공리 I, II가 성립한다. 여기서 고려하는 \(\mF\)의 분수체의 확대체들은 \(\mS\)의 한 원소를 첨가하여 생성할 수 있는 것들뿐이다\footnote{이 확대체들에서 \(\mF\)에 대하여 정수적인 원소의 계는 \(\mR\)에서 \(\mF\)로 넘어간 것과 마찬가지로 \(\mS\)에서 대응하는 함수영역으로 넘어가도 얻어진다. 예컨대 J. König, \emph{Algebraische Größen}, Leipzig 1903, 468/69쪽을 보라.}. 1, 2를 함께 고려하면 수체와 함수체에서 공리 I--V가 성립함을 알았다.

\subsection*{§ 4. 동형정리. 직접합}

이하에서는 환 또는 모듈의 성질만을 가정하고 다른 공리는 전제하지 않는다.

1. \(M\)과 \(\overline M\)이 \(\mR\)에 관한 모듈영역(§ 2)일 때, \(M\)의 각 원소에 \(\overline M\)의 정확히 한 원소가 대응하여 \(\overline M\)을 모두 채우고\footnote{\(\overline M\)에서 통용되는 상등의 정의에 따른 말이다. 각주 9를 보라.}, 이 대응에서 차와 같은 \(\mR\) 원소에 의한 곱셈이 보존되면 \(M\)이 \(\overline M\)에 준동형이라고, 더 정확히는 모듈준동형이라고 하여 \(M\sim\overline M\)으로 쓴다. 즉 \(\beta\sim\overline\beta\), \(\gamma\sim\overline\gamma\)이면 언제나 \((\beta-\gamma)\sim(\overline\beta-\overline\gamma)\), \(r\beta\sim r\overline\beta\)이다.

그러면 \(M\) 안의 \(\mR\)-모듈 \(\mB\)에는 \(\overline M\) 안의 \(\mR\)-모듈 \(\overline{\mB}\)가 준동형으로 대응한다. 또한 \(\overline M\) 안의 한 \(\mR\)-모듈 \(\overline{\mC}\)의 원소들에 대응하는 \(M\)의 모든 원소가 이루는 계 \(\mC^*\)는 \(M\) 안의 \(\mR\)-모듈이고, \(\overline{\mC}\)가 유일하게 결정하는 대응 모듈로서 \(\mC^*\sim\overline{\mC}\)이다. 특히 \(\mA\)가 \(\overline M\)의 영원소에 대응하는 모듈이면 \(\mC^*\)는 \(\mA\)의 약수이고, \(\mA\)를 법으로 서로 합동인 \(M\)의 원소들의 류들로 나뉘며, 이 류들은 \(\overline{\mC}\)의 원소들과 일대일로 대응한다. \(M\)의 \(\mB\)에서 \(\overline M\)의 \(\overline{\mB}\)로 갔다가 다시 \(\mB^*\)로 돌아오면 \(\mB^*=(\mB,\mA)\), 곧 \(\mB\)와 \(\mA\)의 최대공약모듈이 된다. 따라서 \(\mB\)가 \(\mA\)의 약수이면 \(\mB^*=\mB\)이다.

\(M\)과 \(\overline M\)의 원소 대응이 양방향으로 유일하면 두 영역을 동형, 더 정확히는 모듈동형이라 하여 \(M\simeq\overline M\)으로 쓴다.

\(\mA\)가 \(M\) 안의 임의의 \(\mR\)-모듈일 때, \(\mA\)에 따른 합동을 새 상등관계로 삼으면 \(M\)과 준동형인 모듈영역 \(\overline M\), 곧 잉여류모듈 \(M\mid\mA\)가 생긴다. 이때 \(M\)의 각 원소에는 \(\overline M\)에서 그와 같은 모든 원소만이 대응한다. \(\overline M\)에서 상등의 정의로부터 동일성으로 넘어간다는 것은 \(\overline M\)에서 같은 모든 원소를 하나의 류, 곧 잉여류로 묶고 이 잉여류들을 \(\overline M\)의 새 원소로 본다는 뜻이다\footnote{각주 9를 보라.}.

\emph{잉여류모듈로 넘어가면 모든 준동형이 생성된다.} 실제로 \(M\sim\overline M\)이고 \(\mA\)가 \(\overline M\)의 영원소에 대응하는 모듈이면, 위에서 보았듯이 \(\overline M\)은 잉여류모듈 \(M\mid\mA\)와 동형이다.

2. \textbf{제1동형정리.} \emph{\(\overline M\)이 잉여류모듈 \(M\mid\mA\)이고 \(\mC\)가 \(\mA\)의 약수이면
\[
        \overline M\mid\overline{\mC}\simeq M\mid\mC
\]
이다.}
실제로 \(\mA\)에 따른 합동은 \(\mC\)에 따른 합동이기도 하다. 따라서 \(\mA\)에 따라 같은 원소들은 \(\mC\)에 따라서도 같고, 먼저 원소들을 \(\mA\)에 따라 같게 놓아 \(\overline M\)으로 넘어간 뒤 그 가운데 \(\mC\)에 따라 같은 것들을 묶어서 잉여류모듈 \(M\mid\mC\), 곧 \(\mC\)에 따른 상등을 만들 수 있다. 이는 \(\overline M\)에서 \(\overline{\mC}\)에 따라 같게 놓는 것, 곧 \(\overline M\mid\overline{\mC}\)를 만드는 것과 같다.

\textbf{제2동형정리.} \emph{\(\mB,\mA\)가 \(M\)의 모듈이면
\[
        (\mB,\mA)\mid\mA\simeq \mB\mid[\mB,\mA]
\]
이다.}
실제로 1에 따라 \((\mB,\mA)\mid\mA=\overline{\mB}\)로 놓으면 \(\mB\)는 \(\overline{\mB}\)와 준동형이다. 이때 \([\mB,\mA]\)의 모든 원소와 오직 이 원소들만이 \(\overline{\mB}\)의 영원소에 대응하므로, 다시 1에 따라 위 동형을 얻는다.

3. \(\mR,\overline{\mR}\)이 가환환일 때, \(\mR\)의 각 원소에 \(\overline{\mR}\)의 정확히 한 원소가 대응하여 \(\overline{\mR}\)을 모두 채우고 이 대응에서 차와 곱이 보존되면 \(\mR\)이 \(\overline{\mR}\)에 준동형이라고, 더 정확히는 환준동형이라고 하여 \(\mR\sim\overline{\mR}\)으로 쓴다. 원소 대응이 양방향으로 유일하면 두 환을 동형, 더 정확히는 환동형이라 하여 \(\mR\simeq\overline{\mR}\)으로 쓴다.

1, 2의 모든 논의에서 모듈을 \(\mR\)의 아이디얼로 바꾸고\footnote{비가환환에서는 양쪽아이디얼을 바탕에 두어야 한다.}, 모듈준동형과 모듈동형을 환준동형과 환동형으로 바꾸어도 그대로 성립한다. 특히 \(\mc\)가 \(\ma\)의 약수이면, \(\ma\)에 따른 합동을 새 상등관계로 도입하여 잉여류환 \(\mc\mid\ma\)를 얻는다. 이제 잉여류의 곱도 정의된다는 점에 유의해야 한다.

\(\mc\)는 \(\mc\mid\ma\)와 준동형이고, 모든 준동형은 이 방식으로 생성된다. 2에서와 같은 동형정리들이 성립한다.

\textbf{제1동형정리.} \(\overline{\mR}\)이 잉여류환 \(\mR\mid\ma\)이고 \(\mc\)가 \(\ma\)의 약수이면 환동형의 뜻에서
\[
        \overline{\mR}\mid\overline{\mc}\simeq\mR\mid\mc
\]
이다.

\textbf{제2동형정리.} \(\mb,\ma\)가 \(\mR\)의 아이디얼이면 환동형
\[
        (\mb,\ma)\mid\ma\simeq\mb\mid[\mb,\ma]
\]
이 성립한다\footnote{군론에서 알려진 이 동형정리들은 모듈에 대해서는 조금 더 특수한 형태로 데데킨트에게서 처음 나타난다(제3판, 484쪽). 위의 아이디얼 형태는 Masazo Sono, \emph{On Congruences. I, II, III, IV}, Memoirs of the College of Science, Kyoto Imperial University, 2 (1917), 3 (1918, 1919), 특히 2, 215쪽에 있다. 각각 제2동형정리만 명시적으로 진술되어 있다.}.

4. 이제 서로소인 아이디얼들에 관한 계산법칙들을 모은다\footnote{데데킨트에게서 유래한 형태(제4판, § 178, III--VIII)로 제시한다. 근래의 서술에서 늘 나타나는 단위원을 이용한 계산은 훨씬 번거롭다.}. 가환환 \(\mR\)에 곱셈 단위원 \(e\)가 존재한다고 가정한다. \(\mo\)를 단위아이디얼이라 하면 모든 아이디얼 \(\ma\)에 대해 \(\ma=\mo\ma\)이다. 두 아이디얼 \(\ma,\mb\)의 최대공약아이디얼 \((\ma,\mb)\)가 \(\mo\)이면 둘을 서로소라 한다.

\emph{4\(\alpha\). \((\ma,\mb)=\mo\)이면 \((\ma,\mb\mc)=(\ma,\mc)\)이다.} 실제로
\[
(\ma,\mc)=(\ma,\mb)\cdot(\ma,\mc)=(\ma^2,\ma\mb,\ma\mc,\mb\mc)
\]
는 \((\ma,\mb\mc)\)로 나누어지고, 역으로도 그러하다. 이로부터 다음이 따른다.

\emph{4\(\beta\). \(\mc\mb\equiv0(\ma)\), \((\ma,\mb)=\mo\)이면 \(\mc\equiv0(\ma)\)이다.} 실제로 \(\mc\mb\equiv0(\ma)\)는 \((\ma,\mb\mc)=\ma\)와 같은 말이다. 그러므로 \((\ma,\mb)=\mo\)에 따라 \((\ma,\mc)=\ma\), 곧 \(\mc\equiv0(\ma)\)이기도 하다. 더 나아가 다음이 따른다.

\emph{4\(\gamma\). 아이디얼 \(\ma_1,\ldots,\ma_r\) 각각이 아이디얼 \(\mb_1,\ldots,\mb_s\) 각각과 서로소이면, \(\ma_i\)들의 곱은 \(\mb_i\)들의 곱과 서로소이다.} 실제로 \((\ma,\mb)=\mo\), \((\ma,\mc)=\mo\)에서 \((\ma,\mb\mc)=\mo\)가 나오고, 이를 유한번 되풀이하면 주장에 이른다.

4\(\alpha\), 4\(\gamma\)에서 다음을 얻는다.

\emph{4\(\delta\). 아이디얼 \(\mb_1,\ldots,\mb_r\)이 쌍마다 서로소, 곧 \(i\ne k\)일 때 \((\mb_i,\mb_k)=\mo\)이면, 그 최소공배수는 그 곱과 같다.} \((\ma,\mb)=\mo\), \(\mv=[\ma,\mb]\)라 하자. 그러면 \(\mv=\mo\mv=(\ma\mv,\mb\mv)\equiv0(\ma,\mb)\)이고, \(\ma\mb\equiv0(\mv)\)이므로 \(\mv=\ma\mb\)이다. 4\(\gamma\)에 유의하며 이를 유한번 되풀이하면 주장에 이른다.

\emph{4\(\varepsilon\). 아이디얼 \(\mb_1,\ldots,\mb_r\)이 쌍마다 서로소이고
\[
\ma_i=[\mb_1,\ldots,\mb_{i-1},\mb_{i+1},\ldots,\mb_r]
      =\mb_1\cdots\mb_{i-1}\mb_{i+1}\cdots\mb_r
\]
로 놓으면
\[
(\ma_1,\ldots,\ma_{i-1},\ma_{i+1},\ldots,\ma_r)=\mb_i,
\qquad
(\ma_1,\ldots,\ma_r)=\mo
\]
이다.}
실제로 분배법칙에 따라 \((\ma_1,\ma_2)=\mb_3\cdots\mb_r(\mb_2,\mb_1)=\mb_3\cdots\mb_r\)이다. 따라서 \((\ma_1,\ma_2,\ma_3)=\mb_4\cdots\mb_r(\mb_3,\mb_1\mb_2)=\mb_4\cdots\mb_r\)이고, 번호를 알맞게 매겨 이를 유한번 되풀이하면 주장에 이른다.

\(\mm=[\mb_1,\ldots,\mb_r]\)이라는 표현에서 \(\ma_i\)를 \(\mb_i\)의 보충부분이라 한다.

5. 다시 가환환 \(\mR\)에 곱셈 단위원이 존재한다고 가정한다. \(\mR\)의 모든 원소 \(c\)가 \(a_i\in\ma_i\)인 꼴
\[
        c=a_1+a_2+\cdots+a_r
\]
로 정확히 한 가지 방식으로 표현되면, 환 \(\mR\)을 아이디얼 \(\ma_1,\ldots,\ma_r\)의 \emph{직접합}이라 하고 \(\mR=\ma_1+\ma_2+\cdots+\ma_r\)로 쓴다.

\emph{영아이디얼이 쌍마다 서로소인 아이디얼 \(\mb_1,\ldots,\mb_r\)의 최소공배수이고 \(\ma_i\)가 \(\mb_i\)의 보충부분이면, \(\mR\)은 아이디얼 \(\ma_1,\ldots,\ma_r\)의 직접합이다.} 실제로 \(\mo=(\ma_1,\ldots,\ma_r)\)이므로 \(\mR\)의 모든 원소는 \(\ma_i\)들에 의한 덧셈표현을 적어도 하나 가진다. \([\ma_i,\mb_i]=(0)\)이므로 이 표현은 유일하다. 실제로 \(0=a_1+a_2+\cdots+a_r=a_i+b_i\)이면 \(a_i=0\)이다. 제2동형정리에 따라 \(\ma_i\)는 \emph{잉여류환 \(\mR\mid\mb_i\)와 동형}이다. \glqq{}직교관계\grqq{} \(\ma_i\ma_k=(0)\) (\(i\ne k\))와 \(\ma_i^2=\ma_i\)가 성립한다. 마지막 식은 \(\ma_i=\mo\ma_i\)에서 나온다.

\emph{역으로 \(\mR\)의 직접합 표현
\[
        \mR=\ma_1+\ma_2+\cdots+\ma_r
\]
이 존재하고
\[
\mb_i=(\ma_1,\ldots,\ma_{i-1},\ma_{i+1},\ldots,\ma_r)
     =\ma_1+\cdots+\ma_{i-1}+\ma_{i+1}+\cdots+\ma_r
\]
로 놓으면, \(\mb_i\)들은 쌍마다 서로소이고 그 최소공배수는 영아이디얼이다.}
실제로 \(\mo=(\ma_i,\mb_i)\)에서 \(k\ne i\)일 때 \(\mb_k\)가 \(\ma_i\)의 약수이므로 \(\mo=(\mb_k,\mb_i)\)가 따른다. 또 \(c\)가 모든 \(\mb_i\)로 나누어진다고 하자. 그러면
\[
 c=a_2^{(1)}+\cdots+a_r^{(1)}
  =a_1^{(2)}+a_3^{(2)}+\cdots+a_r^{(2)}
  =\cdots
  =a_1^{(r)}+\cdots+a_{r-1}^{(r)},
\]
이고 각 \(a_i^{(\lambda)}\equiv0(\ma_i)\)이다. 표현의 유일성에 따라 매번 한 성분이 0이 되는 데서 모든 \(\lambda\)에 대해 \(0=a_1^{(\lambda)},\ldots,0=a_r^{(\lambda)}\)이고, 따라서 \(c=0\)이다\footnote{5의 정리들은 직접합과 직접교집합의 관계에 관한 일반정리의 특수한 경우이다. H. Prüfer, \emph{Theorie der Abelschen Gruppen I}, Math. Zeitschr. 20 (1924), 165--187쪽, § 6을 보라. 그곳의 논의는 군 개념을 모듈과 아이디얼이 특수한 경우로 들어갈 만큼 일반적으로 잡아도 그대로 유지된다.}.

\begin{center}
{\bfseries\large § 5.\par}
\vspace{0.35em}
{\bfseries\large 소아이디얼과 으뜸아이디얼.\par}
\end{center}

다시 \(\mR\)을 가환환이라 하되, 단위원의 존재를 포함하여 다른 어떤 공리도 가정하지 않는다. 소아이디얼과 으뜸아이디얼에 관한 기본사실들을 모은다\footnote{\emph{아이디얼론} § 4에서는 약수사슬정리라는 공리 I을 가정하였다. 그래서 소아이디얼과 으뜸아이디얼의 어느 성질이 이 공리와 무관한지가 뚜렷하게 드러나지 않는다.}.

1. \(\mR\)의 한 아이디얼 \(\mpideal\)에 따른 잉여류환 \(\mR\mid\mpideal\)이 영인자가 없는 환이면 \(\mpideal\)을 \emph{약한 소아이디얼}이라 한다. 즉 \(a\not\equiv0(\mpideal)\), \(b\not\equiv0(\mpideal)\)이면 언제나 \(ab\not\equiv0(\mpideal)\)이다. \(\mR\)의 한 아이디얼 \(\mpideal^*\)에 따른 잉여류환 \(\mR\mid\mpideal^*\)에 \emph{영인자아이디얼}이 없으면 \(\mpideal^*\)를 \emph{강한 소아이디얼}이라 한다. 즉 \(\ma\not\equiv0(\mpideal^*)\), \(\mb\not\equiv0(\mpideal^*)\)이면 언제나 \(\ma\mb\not\equiv0(\mpideal^*)\)이다.

\(\ma,\mb\)를 주아이디얼로 특수화하면 \emph{모든 강한 소아이디얼은 약한 소아이디얼이기도 하다}. 역으로 \emph{모든 약한 소아이디얼도 강한 소아이디얼}이므로, 단순히 \emph{소아이디얼}이라 말할 수 있다. 실제로 \(\mpideal\)이 약한 소아이디얼이고 \(\ma\not\equiv0(\mpideal)\), \(\mb\not\equiv0(\mpideal)\)이라 하자. \(\ma,\mb\)에서 각각 \(a\not\equiv0(\mpideal)\), \(b\not\equiv0(\mpideal)\)인 원소 \(a,b\)를 택할 수 있다. 그러면 \(ab\not\equiv0(\mpideal)\), 따라서 \(\ma\mb\not\equiv0(\mpideal)\)이므로 \(\mpideal\)은 강한 소아이디얼이기도 하다.

2. 잉여류환 \(\mR\mid\mq\)에서 모든 영인자의 어떤 거듭제곱이 0이 되면 \(\mR\)의 아이디얼 \(\mq\)를 \emph{약한 으뜸아이디얼}이라 한다. 즉 \(a\not\equiv0(\mq)\)이고 모든 \(\varkappa\)에 대해 \(b^\varkappa\not\equiv0(\mq)\)이면 언제나 \(ab\not\equiv0(\mq)\)이다. \(\mR\mid\mq\)에서 영인자가 되는 \(\mR\)의 모든 원소의 계 \(\mpideal\)은 아이디얼을 이루며, \(\mq\)의 약수인 소아이디얼이다. 이를 \emph{대응 소아이디얼}이라 한다. 실제로 \(a\)와 더불어 \(ra\)도 영인자이고, \(a,b\)와 더불어 \(a-b\)도 영인자이다. 마지막 말은 \(a^\varkappa\), \(b^\lambda\)와 더불어 \((a-b)^{\varkappa+\lambda}\)도 언제나 \(\mq\)로 나누어지기 때문이다. 또한 \(a\)가 영인자가 아니면 정의에 따라 \(a\)의 어떤 거듭제곱도 영인자가 아니다. 따라서 \(a\not\equiv0(\mpideal)\), \(b\not\equiv0(\mpideal)\)이면 \(a^\varkappa b^\varkappa=(ab)^\varkappa\not\equiv0(\mq)\), 따라서 \(ab\not\equiv0(\mpideal)\)이다.

잉여류환 \(\mR\mid\mq^*\)에서 모든 영인자아이디얼의 어떤 거듭제곱이 0이 되면 \(\mR\)의 아이디얼 \(\mq^*\)를 \emph{강한 으뜸아이디얼}이라 한다. 즉 \(\ma\not\equiv0(\mq^*)\)이고 모든 \(\varkappa\)에 대해 \(\mb^\varkappa\not\equiv0(\mq^*)\)이면 언제나 \(\ma\mb\not\equiv0(\mq^*)\)이다. 약한 으뜸아이디얼의 경우와 같이 다음을 보인다. \(\mR\mid\mq^*\)에서 영인자아이디얼이 되는 \(\mR\)의 모든 아이디얼의 최대공약아이디얼은 \(\mq^*\)의 약수인 한 소아이디얼, 곧 \emph{대응 소아이디얼}을 이룬다. 역으로 같은 대응 소아이디얼 \(\mpideal\)을 갖는 모든 으뜸아이디얼을 \glqq{}\(\mpideal\)에 대응한다\grqq{}고 한다.

\(\ma,\mb\)를 주아이디얼로 특수화하면 \emph{모든 강한 으뜸아이디얼은 약한 으뜸아이디얼이기도 하다}. \emph{그러나 일반적으로 그 역은 성립하지 않는다}\footnote{R. Hölzer가 알려 준 다음 예가 이를 보인다. \(\mR\)을 한 체를 계수로 하는 가산개의 부정원 \(x_i\)의 다항식영역이라 하고
\[
\mq=(x_1^2,x_2^3,\ldots,x_\nu^{\nu+1},\ldots,x_i x_k\,[i\ne k],\ldots)
\]
라 하자. 잉여류환의 영인자들은 정확히 \(\mpideal=(x_1,x_2,\ldots,x_\nu,\ldots)\)로 나누어지는 다항식들이며, 이 영인자들의 어떤 거듭제곱은 각각 \(\mq\)로 나누어진다. 따라서 \(\mq\)는 \(\mpideal\)을 대응 소아이디얼로 갖는 약한 으뜸아이디얼이다. 그러나 \(\ma=(x_1,x_3,x_5,\ldots,x_{2\nu+1},\ldots)\), \(\mb=(x_2,x_4,\ldots,x_{2\nu},\ldots)\)로 놓으면 \(\ma\mb\equiv0(\mq)\)이지만 \(\ma,\mb\)의 어떤 거듭제곱도 \(\mq\)로 나누어지지 않으므로 \(\mq\)는 강한 으뜸아이디얼이 아니다.}. \emph{그러나 \(\mR\)에 약수사슬정리라는 공리 I을 가정하면 모든 약한 으뜸아이디얼은 강한 으뜸아이디얼이기도 하다.} 따라서 단순히 \emph{으뜸아이디얼}이라 말할 수 있다. 실제로 \(\mq\)가 약한 으뜸아이디얼이고 \(\ma\not\equiv0(\mq)\), 또 모든 \(\varkappa\)에 대해 \(\mb^\varkappa\not\equiv0(\mq)\)이라 하자. 그러면 \(a\not\equiv0(\mq)\)인 \(a\in\ma\)와 모든 \(\varkappa\)에 대해 \(b^\varkappa\not\equiv0(\mq)\)인 \(b\in\mb\)가 존재한다. 따라서 \(ab\not\equiv0(\mq)\), 곧 \(\ma\mb\not\equiv0(\mq)\)이다. 실제로 \(\mb\)의 각 \(b\)마다 \(b^\lambda\equiv0(\mq)\)이게 하는 지수 \(\lambda\)가 존재한다면, 유한개의 기저원소의 지수를 \(\lambda_1,\ldots,\lambda_r\)라 할 때 \(\mb^{\lambda_1+\cdots+\lambda_r}\equiv0(\mq)\)일 것이다\footnote{공리 I로부터 유한 아이디얼기저를 끌어내려면 \(\mR\)에 고정된 정렬순서를 바탕에 두어야 한다(§ 6 참조).}. 이 마지막 논의는 다음도 보인다. \(\mpideal\)이 대응 소아이디얼이면 \(\mpideal^\varrho\equiv0(\mq)\)이게 하는 최소 지수 \(\varrho\)가 존재하며, 이를 \(\mq\)의 \emph{지수}라 한다.

3. 이하에서는 다시 약수사슬정리를 \emph{가정하지 않는다}. 유한개의 아이디얼의 최소공배수 \([\ma_1,\ldots,\ma_r]\)에서 어떤 \(\ma_i\)도 나머지 아이디얼들의 최소공배수로 나누어지지 않아 어느 \(\ma_i\)도 생략할 수 없으면 이를 \emph{최단표현}이라 한다.

\emph{같은 소아이디얼 \(\mpideal\)에 대응하는 유한개의 약한 또는 강한 으뜸아이디얼의 최소공배수는 다시 \(\mpideal\)을 대응 소아이디얼로 갖는 약한 또는 강한 으뜸아이디얼이다. 서로 다른 소아이디얼들에 대응하는 유한개의 약한 또는 강한 으뜸아이디얼에 의한 최단표현은 약한 또는 강한 으뜸아이디얼이 아니다.}

\(\mf=[\mq_1,\ldots,\mq_r]\)이고 \(\mpideal\)이 약한 으뜸아이디얼 \(\mq_i\)들의 대응 소아이디얼이라 하자. 그러면 \(\mpideal\)은 어떤 거듭제곱이 \(\mf\)로 나누어지는 모든 원소의 전체이기도 하다. \(a\not\equiv0(\mf)\)이고 모든 \(\varkappa\)에 대해 \(b^\varkappa\not\equiv0(\mf)\)이면 \(b\not\equiv0(\mpideal)\)이고, 적어도 한 지표 \(i\)에 대해 \(a\not\equiv0(\mq_i)\)이다. 따라서 \(ab\not\equiv0(\mq_i)\), 곧 \(ab\not\equiv0(\mf)\)이다. 줄곧 원소를 아이디얼로 바꾸면 더는 \(\mb\not\equiv0(\mpideal)\)을 결론낼 수 없다.

이제 \(i\ne k\)일 때 \(\mpideal_i\ne\mpideal_k\)인 약한 으뜸아이디얼들의 최단표현
\[
        \mm=[\mq_1,\ldots,\mq_r],\qquad r\ge2
\]
를 생각하자. 대응 소아이디얼들 가운데 적어도 하나, 이를테면 \(\mpideal_1\)은 나머지 어느 것에도 들어가지 않는다. \(\mpideal\)들의 모든 진배수사슬은 늦어도 \(r\)단계에서 멈추기 때문이다. 따라서 \(a_i^{\varrho_i}\equiv0(\mq_i)\)이고 \(a_2\not\equiv0(\mpideal_1),\ldots,a_r\not\equiv0(\mpideal_1)\)이게 하는 \(a_i\in\mpideal_i\)를 택할 수 있다. 또 \(q_1\not\equiv0(\mm)\)인 \(q_1\in\mq_1\)을 택하면
\[
q_1a_2^{\varrho_2}\cdots a_r^{\varrho_r}\equiv0(\mm)
\]
이지만 \((a_2^{\varrho_2}\cdots a_r^{\varrho_r})\)의 어떤 거듭제곱도 \(\mm\)으로 나누어지지 않는다. 그렇지 않으면 \(\mpideal_1\)로 나누어져야 하기 때문이다. 따라서 \(\mm\)은 약한 으뜸아이디얼이 아니다. 줄곧 원소를 아이디얼로, 즉 \(q_1\)을 \(\mq_1\)로, \(a_i\)를 \(\ma_i\not\equiv0(\mpideal_1)\)이면서 \(\ma_i^{\varrho_i}\equiv0(\mq_i)\)인 \(\ma_i\)로 바꾸면 강한 으뜸아이디얼에 대한 증명이 나온다\footnote{약수사슬정리를 불필요하게 만든 이 원래 증명의 수정은 B. L. van der Waerden에게서 배웠다.}.

4. \emph{모듈 몫과 아이디얼 몫.} \(\mT\)를 \(\mR\)의 확대환이라 하고, \(\mA,\mB,\mC,\ldots\)를 \(\mT\) 안의 \(\mR\)-모듈이라 하자. \emph{몫} \(\mC=\mA:\mB\)는 다음 조건을 만족하는 모듈로 정의한다. \(\mC\mB\equiv0(\mA)\)이고, \(\mD\mB\equiv0(\mA)\)이면 \(\mD\equiv0(\mC)\)이다. 따라서 \(\mC\)는 \(\mC\mB\equiv0(\mA)\)이게 하는 \glqq{}가장 작은\grqq{} 모듈이다. 이 조건으로 몫은 \(\mD\mB\equiv0(\mA)\)인 모든 모듈 \(\mD\)의 최대공약모듈로서 유일하게 결정된다. 영모듈은 확실히 그러한 모듈이므로 이 \(\mD\)들은 존재한다.

정의로부터 다음이 따른다. \(\mB\)가 \(\overline{\mB}\)의 배수이면 \(\mA:\overline{\mB}\)는 \(\mA:\mB\)의 배수이다. \(\mA\)가 \(\overline{\mA}\)의 배수이면 \(\mA:\mB\)는 \(\overline{\mA}:\mB\)의 배수이다. 또한
\[
        \mA:\mB\mC=(\mA:\mB):\mC=(\mA:\mC):\mB
\]
이다. 확대환 \(\mT\)가 \(\mR\)과 같으면 \emph{모듈 몫}은 \emph{아이디얼 몫} \(\ma:\mb\)이 되고, 위 계산법칙들이 그대로 성립한다. \(\ma\mb\equiv0(\ma)\)이므로 이 경우에는 \(\ma\equiv0(\ma:\mb)\)도 따른다.

5. \emph{\(\ma:\mb=\ma\)이면 \(\mb\)가 \(\ma\)에 대하여 상대소라고 한다.} 따라서 \(\mb\)가 \(\ma\)에 대하여 상대소일 필요충분조건은 \(\md\mb\equiv0(\ma)\)이면 언제나 \(\md\equiv0(\ma)\)인 것이다. 4의 계산법칙에서 다음이 따른다. \emph{\(\mb,\mc\)가 \(\ma\)에 대하여 상대소이면 그 곱과 최소공배수도 \(\ma\)에 대하여 상대소이다.} \(\mb\)가 \(\ma\)에 대하여 상대소이고 \(\ma\)도 \(\mb\)에 대하여 상대소이면 \(\ma,\mb\)를 서로 상대소라고 한다. § 4, 4\(\beta\)에서 서로소인 아이디얼들은 서로 상대소이기도 하다는 것이 따른다.

\setcounter{footnote}{25}

\begin{center}
{\bfseries\large § 6.\par}
\vspace{0.35em}
{\bfseries\large 약수사슬정리를 가정한 아이디얼론.\par}
\end{center}

약수사슬정리라는 \emph{공리 I이 성립}하는 가환환 \(\mR\)을 바탕에 두고, 이하에서는 언제나 \(\mR\)의 \emph{모든 원소에 대한 고정된 정렬순서}도 바탕에 둔다. 그러면 \(\mR\)의 \emph{모든 아이디얼에도 정렬순서}가 주어진다. 실제로 이 두 가정에서 \(\mR\)의 모든 아이디얼이 유한개의 원소로 이루어진 아이디얼기저를 갖는다는 것이 곧바로 나온다. 이제 \(\mR\) 원소의 정렬순서에서 \(\mR\)의 모든 유한 부분집합의 준사전식 순서로 넘어가고\footnote{그 뜻은 다음과 같다. 유한 부분집합 \(\mathfrak A_n\)의 원소 \(a_1,\ldots,a_n\)을 지표 순서가 \(\mR\)의 정렬순서와 일치하도록 나타낸다. 원소가 \(n\)개인 모든 부분집합 \(\mathfrak A_n\)은 \(m>n\)개인 모든 \(\mathfrak A_m\)보다 앞선다. 원소 수가 같은 \(\mathfrak A_n,\mathfrak B_n\)에 대해서는 대응하는 \(b_i\)와 다른 첫 원소 \(a_i\)가 \(\mR\)의 정렬순서에서 \(b_i\)보다 앞서면 \(\mathfrak A_n\)이 \(\mathfrak B_n\)보다 앞선다고 한다. 이로써 유한 부분집합들의 모든 계에는 첫 원소 \(\mathfrak A_n\)이 대응한다. 따라서 \(\mR\)의 정렬순서가 주어지면 다른 선택공리는 필요 없다. 특히 대수적 수체의 경우처럼 \(\mR\)이 가산집합이면 선택공리를 전혀 바탕에 두지 않는다. 아이디얼론에서 선택공리가 하는 역할은 \emph{아이디얼론}, 각주 9에서 자세한 설명 없이 지적하였다.}, 각 아이디얼에는 가능한 여러 기저 가운데 유한 부분집합의 정렬순서에서 첫째인 것을 특수기저로 대응시키면, 유한 부분집합들과 더불어 아이디얼들도 정렬된다.

\textbf{정리 I.} \emph{약수사슬정리를 가정하면 \(\mR\)의 모든 아이디얼은 유한개의 기약아이디얼의 최소공배수로 표현된다. 여기서 기약아이디얼이란 두 진약수의 최소공배수로 표현되지 않는 아이디얼이다.}

증명을 위해 다음을 보이면 된다. 한 아이디얼 \(\mm\)에 정리 I이 성립하지 않으면, \(\mm\)에는 정리 I이 역시 성립하지 않는 진약수가 있다. 그러면 가정한 약수사슬정리에 어긋나게 유한 단계에서 멈추지 않는 약수사슬을 구성할 수 있다. 실제로 \(\mm\)은 가약이어야 한다. 그렇지 않으면 \(\mm=[\mm]\)이 요구되는 표현을 주기 때문이다. \(\mm=[\ma,\mb]\)이면 정리 I은 \(\ma,\mb\)에 동시에 성립할 수 없다. 그렇다면 \(\mm\)에도 대응하는 표현이 생기기 때문이다. 따라서 \(\mm\)에는 정리 I이 성립하지 않는 진약수들이 있고, 그 가운데 아이디얼의 정렬순서에서 첫째를 \(\ma_1\)이라 하자. 같은 방식으로 \(\ma_1\)의 진약수 \(\ma_2\), 일반적으로 \(\ma_i\)의 진약수 \(\ma_{i+1}\)를 구성하면, 잘 정의되면서 유한 단계에서 멈추지 않는 약수사슬을 얻어 가정에 모순된다\footnote{모든 모듈계에 약수사슬정리를 가정하면 정리 I은 모듈영역에도 그대로 성립한다. 정리 I은 순전히 집합론적 성격을 가지며, 아이디얼의 정렬순서를 증명에 쓰는 유일한 정리이기도 하다. 모든 연산과 무관하게 다음 집합론적 개념들만 문제된다. 한 집합 \(\mathfrak M\)에서 그 멱집합의 한 부분집합 \(\Sigma\), 곧 부분집합들의 한 계를 택한다. \(\Sigma\)가 정렬되어 있다고 하고, 그 원소를 \(\Sigma\)-집합이라 하자. \(\Sigma\)에 다음 사슬정리를 가정한다. \(A_\nu\)가 \(A_{\nu-1}\)의 진상위집합인 \(\Sigma\)-집합들의 모든 사슬 \(A_1,A_2,\ldots,A_\nu,\ldots\)는 유한 단계에서 멈춘다. 한 \(\Sigma\)-집합 \(A\)가 모두 \(A\)의 진상위집합인 두 \(\Sigma\)-집합의 교집합이면 가약, 그렇지 않으면 기약이라 한다. 그러면 위 논의에 따라 모든 \(\Sigma\)-집합은 유한개의 기약 \(\Sigma\)-집합의 교집합으로 표현된다.}.

가정한 약수사슬정리에 따라 § 5, 2에서 단순히 으뜸아이디얼이라 말할 수 있다. 으뜸성과 기약성의 관계는 다음과 같다.

\textbf{정리 II.} \emph{약수사슬정리를 가정하면 모든 기약아이디얼은 으뜸이다. 다시 말해, 으뜸이 아닌 모든 아이디얼은 가약이다.}

잉여류환 \(\mR\mid\mm\)으로 넘어가도 준동형 때문에 약수사슬정리가 유지되고, \(\mm\)을 최소공배수로 나타내는 표현에는 영아이디얼의 표현이 대응하며, 제1동형정리(§ 4, 3)에 따라 \(\mm\)의 으뜸 약수에는 \(\mR\mid\mm\)의 으뜸아이디얼이 대응하고 그 역도 성립한다. 따라서 잉여류환에서 영아이디얼을 분해하는 문제로 제한할 수 있다. 이 환도 같은 일반성을 가지므로 처음부터 분해할 아이디얼을 \(\mR\)의 영아이디얼이라 해도 된다.

이제 \(\mR\)의 영아이디얼이 으뜸이 아니라고 하자. 그러면 \(\ma\ne(0)\), 모든 \(\varkappa\)에 대해 \(\mb^\varkappa\ne(0)\)이지만 \(\ma\mb=(0)\)인 아이디얼쌍 \(\ma,\mb\)가 적어도 하나 존재하며, \(\mb\)는 주아이디얼이라고 가정해도 된다. 가정에 따라 유한 단계에서 멈추는 아이디얼 몫의 약수사슬
\[
        \ma,\ \ma:\mb,\ldots,\ma:\mb^\nu,\ldots
\]
을 만든다. 예컨대 \(\mt=\ma:\mb^{m-1}=\ma:\mb^m=\cdots\)가 이 사슬에서 그 뒤의 모든 아이디얼과 같다고 하자. 그러면 \(\mt=\mt:\mb\), 곧 \(\mb\)가 \(\mt\)에 대하여 상대소이다. 또한 \(\mt\)는 \(\ma\)의 약수이므로 영아이디얼과 다르고, 모든 지수 \(\varkappa\)에 대해 \(\mb^\varkappa\)도 영아이디얼과 다르다. 따라서 정리 II를 증명하려면
\[
        (0)=[\mt,\mb^{m+1}]
\]
을 보이면 충분하다. 정의에 따라
\[
        \mb^{m-1}\mt\equiv0(\ma),
        \qquad\text{따라서}\qquad
        \mb^m\mt=(0).
\]
그러므로
\[
        \mc=[\mt,\mb^{m+1}]\equiv0(\mb^m\mt)
\]
만 보이면 된다. 이는 \(\mb\)가 주아이디얼이고 \(\mt\)에 대하여 상대소라는 가정에서 따른다. \(c\in\mc\)는 \(\mb^{m+1}\)로 나누어지므로, \(n\)을 정수의 기호라 할 때
\[
        c=k b^{m+1}+n b^{m+1}=r b^m
\]
으로 표현된다. 여기서 \(r\)은 다시 \(\mR\)의 원소이다. \(c\equiv0(\mt)\)에서 \(r b^m\equiv0(\mt)\)이고, \(\mt:\mb=\mt\)이므로 \(r\equiv0(\mt)\)이다. 따라서 \(c\equiv0(\mt\mb^m)\)이고 정리 II가 증명되었다.

\textbf{정리 III.} \emph{약수사슬정리를 가정하면 모든 아이디얼은 서로 다른 소아이디얼들에 대응하는 유한개의 으뜸아이디얼, 곧 극대 으뜸성분들의 최소공배수로 된 최단표현을 갖는다\footnote{대응 소아이디얼들과 고립성분들은 유일하게 결정된다. \emph{아이디얼론}, 또는 W. Krull, \emph{Ein neuer Beweis für die Hauptsätze der allgemeinen Idealtheorie}, Math. Ann. 90 (1923), 55--64쪽을 보라. § 7에서는 그곳에 나타나는 단순한 특수한 경우에 유일성을 직접 증명한다. 그곳에서 유일하다고 판정되는 으뜸아이디얼들은 고립성분이다.}.}

그러한 표현을 얻으려면 정리 I에 따라 존재하는 기약아이디얼, 따라서 정리 II에 따라 으뜸아이디얼에 의한 표현을 가능한 대로 최단표현으로 바꾸기만 하면 된다. 같은 소아이디얼에 대응하는 으뜸아이디얼들을 묶으면 § 5, 3에 따라 요구되는 표현을 얻는다. 이 성분들 가운데 어떤 것들의 최소공배수도 더는 으뜸이 아니므로 으뜸성분들을 극대라고 부를 수 있다.

\begin{center}
{\bfseries\large § 7.\par}
\vspace{0.35em}
{\bfseries\large 이중사슬정리를 가정한 아이디얼론.\par}
\end{center}

지금까지 전개한 이론에 비해 생기는 단순화는 1, 2에서 제시할 보조정리들에 기초한다.

1. \emph{영인자가 없는 가환환에서 배수사슬정리가 성립하면 그 환은 체이기도 하다.} \(a\ne0\)일 때 방정식 \(ax=b\)가 언제나 환 안에서 해를 갖는다는 것을 보이면 된다. 해가 둘 이상일 수 없다는 것은 환에 영인자가 없다는 가정에서 따른다.

\(\ma\)를 \(a\)에서 생성되는 주아이디얼이라 하자. 가정에 따라 거듭제곱의 열은 유한 단계에서 멈춘다. 예컨대 \(\ma^m\)이 그 뒤의 모든 것과 같다고 하자. \(\mb\)를 \(b\)에서 생성되는 주아이디얼이라 하면
\[
\ma^m\mb=\ma^{m+1}\mb,
\qquad
\ma^{m+1}\mb=(a^{m+1}b).
\]
특히 원소 \(a^m b\)에 대해, \(n\)을 정수의 기호라 할 때
\[
        a^m b=k a^{m+1}b+n a^{m+1}b=a^{m+1}c
\]
인 표현이 나온다. 여기서 \(c\)는 다시 \(\mR\)의 원소이다. 가정에 따라 영인자가 없으므로 \(b=ac\)이고, 체의 성질이 증명되었다.

2. 1에서 곧바로 다음 보조정리를 얻는다. \emph{가환환에서 배수사슬정리가 성립하면 한 소아이디얼은 \(\mo\)와 다른 진약수를 갖지 않는다.}

마찬가지로 다음 덧붙임이 따른다. \emph{비영아이디얼을 법으로 한 배수사슬정리라는 공리 II만 성립하면, 모든 비영 소아이디얼은 \(\mo\)와 다른 진약수를 갖지 않는다.}

실제로 소아이디얼 \(\mpideal\)에 따른 잉여류환은 정의에 따라 영인자가 없는 환이고, 준동형 때문에 여기서도 배수사슬정리가 성립하므로 1에 따라 체이다. \(\mpideal\)의 약수들과 잉여류환의 아이디얼들이 일대일로 대응하므로 \(\mpideal\)은 \(\mo\)와 다른 진약수를 갖지 않는다\footnote{\(\mR\), 따라서 \(\mo\), 곧 \(\mR\)의 모든 원소로 이루어진 단위아이디얼에는 단위원이 존재하지 않아도 된다. 1에 따라 잉여류환에는 한 단위류가 존재하지만 그렇다. 이 종류의 가장 단순한 예는 --- 덧붙임의 더 약한 경우에 해당한다 --- 모든 짝수의 계이다.}.

\textbf{정리 IV.} \emph{가환환에서 이중사슬정리가 성립하면, 한 아이디얼의 모든 최단표현에서 단위아이디얼 \(\mo\)에 대응하지 않는 으뜸성분들은 유일하게 결정된다.}

\emph{덧붙임.} 공리 I, II를 가정하면 정리 IV는 모든 비영아이디얼에 성립한다.

\[
\mm=[\mq,\mq_1,\ldots,\mq_r]
   =[\overline{\mq},\overline{\mq}_1,\ldots,\overline{\mq}_{\overline r}]
\]
가 \(\mm\)의 극대 으뜸성분들에 의한 최단표현이고, \(\mo,\mpideal_1,\ldots,\mpideal_r\) 및 \(\mo,\overline{\mpideal}_1,\ldots,\overline{\mpideal}_{\overline r}\)가 각각 대응 소아이디얼이라 하자. \(\mo\)에 대응하는 으뜸아이디얼이 나타나지 않으면 \(\mq\), 또는 \(\overline{\mq}\)는 생략한다. 식
\[
        \mo^\varrho\mpideal_1^{\varrho_1}\cdots\mpideal_r^{\varrho_r}
        \equiv0(\mm)
\]
때문에 각 \(\overline{\mpideal}_i\)는 어떤 \(\mpideal_i\)에 들어가야 하고, 보조정리에 따라 그 아이디얼과 같다. 똑같이 각 \(\mpideal_k\)도 어떤 \(\overline{\mpideal}_k\)와 같아야 하므로 \(\mo\)와 다른 소아이디얼들은 일치하여 \(\mpideal_\lambda=\overline{\mpideal}_\lambda\)이다. 또한 \(\mq\mq_1\cdots\mq_r\equiv0(\overline{\mq}_\lambda)\)에서 \(\mq_\lambda\equiv0(\overline{\mq}_\lambda)\)가 나온다. 나머지 성분들은 \(\mpideal_\lambda\)로 나누어지지 않으므로 \(\overline{\mq}_\lambda\)에 대하여 상대소이기 때문이다. 마찬가지로 \(\overline{\mq}_\lambda\equiv0(\mq_\lambda)\)이므로 \(\mo\)에 대응하지 않는 으뜸성분들의 유일성이 증명되었다. 끝으로 \(\mq\)가 실제로 나타나면 최단표현 때문에 \(\overline{\mq}\)도 실제로 나타나야 하므로 대응 소아이디얼의 유일성도 증명된다. 이 유일성 증명은 \(\mo\)에 대응하는 으뜸성분을 포함하지 않는 모든 표현이 동시에 최단임도 보인다.

3. 이중사슬정리의 가정에 단위원의 존재를 더하면 정리 IV는 다음 정리 V로 강화된다.

\textbf{정리 V.} \emph{단위원을 갖는 가환환에서 이중사슬정리가 성립하면, 모든 아이디얼은 쌍마다 서로소인 유한개의 으뜸아이디얼의 곱으로 유일하게 표현된다.}

\emph{덧붙임.} 공리 I, II, III을 가정하면 정리 V는 모든 비영아이디얼에 성립한다.

단위원이 존재하므로 \(\mo^2=\mo\)이고, 따라서 \(\mo\)에 대응하는 모든 으뜸아이디얼은 \(\mo\)와 같다. 그러한 성분은 \(\mo\)와 다른 아이디얼의 최단표현에는 나타날 수 없다. 따라서 정리 IV에 따라 으뜸성분들이 유일하게 결정된다. 동시에 \(\mo\)를 생략하면 모든 표현이 최단표현이 된다. 또한 2의 보조정리에 따라 서로 다른 두 소아이디얼은 언제나 서로소이다. 따라서 § 4, 4의 계산법칙에 따라 으뜸성분들도 서로소이고, 최소공배수는 그 곱이 된다.

가정들로부터 다음 보조정리도 따른다.

\textbf{보조정리.} \emph{단위원과 이중사슬정리를 갖는 가환환에서, 단위아이디얼을 제외하면 한 아이디얼 \(\mm\)에 들어가지 않는 소아이디얼들이 정확히 \(\mm\)에 대하여 상대소인 소아이디얼들이다. 서로 상대소와 서로소라는 개념은 일치한다.}

\emph{덧붙임.} 공리 I, II, III을 가정하는 경우에는 \(\mm\)이 비영아이디얼이라고 해야 한다.

\(\mpideal\)이 \(\mm\)에 들어가지 않으면 \(\mm\)의 모든 대응 소아이디얼과 다르다. 따라서 2의 보조정리에 따라 이들 \(\mo\)와 다른 소아이디얼 가운데 어느 것으로도 나누어지지 않는다. 그러므로 \(\mpideal\)은 모든 으뜸성분에 대하여 상대소이고, 따라서 \(\mm\)에 대하여 상대소이다. 동시에 모든 으뜸성분과 서로소이고, 따라서 \(\mm\)과도 서로소이다.

\(\mpideal\ne\mo\)가 \(\mm\)에 들어가면 어떤 대응 소아이디얼과 같아야 한다. 이를테면 성분 \(\mq=\mq_1\)에 대응한다고 하자. 그러면 \(\mpideal^{\varrho-1}\equiv0(\mq:\mpideal)\)이지만 \(\mpideal^{\varrho-1}\not\equiv0(\mq)\)이므로 \(\mq:\mpideal\)은 \(\mq\)의 진약수이다. 동시에 \(\mq:\mpideal\)은 \(\mpideal^{\varrho-1}\)의 약수로서 \(\mo\)와 다른 소아이디얼 하나로만 나누어지므로 으뜸이다. 유일분해에 따라 \(\mm:\mpideal\)로 나누어지는 곱 \((\mq:\mpideal)\cdot\mq_2\cdots\mq_r\)도 \(\mm\)의 진약수이다. 따라서 \(\mpideal\)은 \(\mm\)에 대하여 상대소가 아니다.

그러므로 아이디얼 \(\mt\ne\mo\)가 \(\mm\)에 대하여 상대소일 필요충분조건은 \(\mt\)의 대응 소아이디얼 가운데 어느 것도 \(\mm\)에 들어가지 않는 것이다. 이 경우 \(\mt\)는 \(\mm\)과 서로소이기도 하다. 역으로 서로소인 두 아이디얼은 언제나 서로 상대소이므로(§ 5, 5), 위 가정 아래에서는 두 개념이 일치한다.

\begin{center}
{\bfseries\large § 8.\par}
\vspace{0.35em}
{\bfseries\large 분수체 안에서 정수적으로 닫힌 경우의 아이디얼론.\par}
\end{center}

지금까지 전개한 아이디얼론에 마지막 두 공리를 더하여 통상적인 더 강한 이론으로 만든다.

1. \emph{보조정리.} \(\mR\)이 단위원을 갖고 영인자가 없는 가환환이며 \(\mR\)에서 약수사슬정리를 가정하면(공리 I, III, IV), \(\ma=\ma\mb\), \(\ma\ne(0)\)에서 언제나 \(\mb=\mo\)가 따른다. 따라서 영아이디얼 및 단위아이디얼과 다른 모든 아이디얼에 대해
\[
        \mc^\varrho\ne\mc^{\varrho+1}
\]
이다\footnote{반대로 같은 가정 아래에서도 \(\ma\mb=\ma\mc\)에서 \(\mb=\mc\)를 결론낼 수는 없다. 다음 예가 이를 보인다. \(\mR\)을 한 체를 계수로 하는 \(x,y\)의 다항식영역이라 하고
\[
\ma=(x,y),\qquad \mb=(x^2,xy,y^2),\qquad \mc=(x^2,y^2)
\]
라 하자. 실제로 \(\ma\mb=\ma\mc=\ma^3\)이지만 \(\mb\ne\mc\)이다.}.

증명은 \(\ma\mb\)를 \(\mb\)에 관한 유한 모듈로 보면 § 1, 3의 보조정리와 같이 나온다. 공리 I에 따라 존재하는 \(\ma\)의 아이디얼기저를 \(\alpha_1,\ldots,\alpha_n\)이라 하면, \(\ma=\ma\mb\)에서 연립방정식
\[
        \alpha_i=b_{i1}\alpha_1+\cdots+b_{in}\alpha_n,
        \qquad b_{ik}\equiv0(\mb),\quad i=1,\ldots,n
\]
을 얻는다. \(\mR\)에 영인자가 없으므로 여기서 \(\lvert b_{ik}-e_{ik}\rvert=0\)이 따른다. \(i\ne k\)이면 \(e_{ik}=0\), \(e_{ii}=e\)이다. 그런데
\[
        e^n+b_1e^{n-1}+\cdots+b_n=0,
        \qquad b_i\equiv0(\mb)
\]
에서 \(e\equiv0(\mb)\), 따라서 \(\mb=\mo\)가 따른다.

2. 이제 공리 I--IV에 분수체 안에서 정수적으로 닫혀 있다는 가정(공리 V)을 더하면 모든 비영 으뜸아이디얼이 그 대응 소아이디얼의 한 거듭제곱이 됨을 보이기만 하면 된다. 그러면 모든 비영아이디얼에 대한 표현
\[
        \mm=\mpideal_1^{\varrho_1}\cdots\mpideal_r^{\varrho_r}
\]
의 유일성은 정리 V와 1의 보조정리로 이미 증명되었다. 동시에 이 소아이디얼들은 단위아이디얼과 다른 진약수를 갖지 않는다는 것도 보였다. 지수가 2인 경우는 데데킨트의 따름정리 II(§ 1, 4)를 이용하여 증명하고, 일반적인 경우는 완전귀납법으로 증명한다.

\emph{보조정리.} \emph{공리 I--V를 가정하면 지수가 2인 으뜸아이디얼은 \(\mq=\mpideal^2\)뿐이다.}

이를 증명하려면
\[
        \mpideal^2\equiv0(\mq),\qquad
        \mq\equiv0(\mpideal),\qquad
        \mq\not\equiv0(\mpideal^2)
\]
에서 반드시 \(\mq=\mpideal\)임을 보이면 된다. 따라서
\[
        c\equiv0(\mq),\qquad
        c\equiv0(\mpideal),\qquad
        c\not\equiv0(\mpideal^2)
\]
라 하자. § 7, 3의 보조정리에서 \(\mo c:\mpideal\)이 \(\mo c\)의 진약수임이 따른다. 따라서
\[
        b\mpideal\equiv0(\mo c)\equiv0(\mq),
        \qquad b\not\equiv0(\mo c)
\]
이게 하는 원소 \(b\)가 존재한다. 그러므로 \(\gamma=b/c\)는 정수적이지 않다. 즉 분수체에는 속하지만 \(\mR\)에는 속하지 않는다.

\(b\not\equiv0(\mpideal)\)임을 보이면 된다. 그러면 \(\mq\)가 \(\mpideal\)에 대응하는 으뜸아이디얼이므로 \(\mpideal\equiv0(\mq)\)가 따른다.

데데킨트의 따름정리 II에 따라 \(\mR\)의 원소 \(m,n\)을 택하여
\[
        \gamma=b/c=m/n
\]
이고 \(m^2/n\)도 정수적이지 않게 할 수 있다. 따라서
\[
        bn=mc,\qquad
        m\not\equiv0(\mo n),\qquad
        m^2\not\equiv0(\mo n).
\]
\(b\mpideal\)에 \(m\)을 곱하면
\[
        mb\mpideal\equiv0(\mo mc),
        \qquad\text{따라서}\qquad
        mb\mpideal\equiv0(\mo nb),
        \qquad\text{그러므로}\qquad
        m\mpideal\equiv0(\mo n).
\]
마지막 결론은 환에 영인자가 없기 때문에 성립한다. 이로부터
\[
        m\not\equiv0(\mpideal)
\]
이 따른다. 실제로 \(m^2\not\equiv0(\mo n)\)이고 \(n\equiv0(\mpideal)\)이다. 마지막 관계는 § 7, 3의 보조정리에서 나온다. \(m\not\equiv0(\mo n)\)이므로 \(\mo n:\mpideal\)은 \(\mo n\)의 진약수이기 때문이다. 이제 네 관계
\[
        m\not\equiv0(\mpideal),\qquad
        n\equiv0(\mpideal),\qquad
        c\equiv0(\mpideal),\qquad
        c\not\equiv0(\mpideal^2),\qquad
        bn=mc
\]
에서 \(b\not\equiv0(\mpideal)\)이 나온다. 실제로 \(\mpideal^2\)가 으뜸이므로 우선 \(mc\not\equiv0(\mpideal^2)\)이고, \(bn=mc\)에서 \(b\not\equiv0(\mpideal)\)이 따른다. 이로써 보조정리가 증명되었다.

3. \emph{보조정리.} \emph{공리 I--V를 가정하면 지수가 \(\varrho\)인 으뜸아이디얼은 \(\mq=\mpideal^\varrho\)뿐이다.}

이 보조정리는 공리를 더 쓰지 않고 2의 보조정리에서 나온다\footnote{보조정리 2가 성립한다고 가정한 보조정리 3의 증명은 Masazo Sono, \emph{On Congruences II}, §§ 9, 10에 있다. 각주 19도 보라.}. 먼저 다음을 보여야 한다.
\[
\text{\(c\equiv0(\mpideal)\), \(c\not\equiv0(\mpideal^2)\)이면 모든 \(\sigma\)에 대해}\qquad
        \mpideal^\sigma=(\mo c^\sigma,\mpideal^{\sigma+1}).
\]
2의 보조정리에 따라 \(\mpideal=(\mo c,\mpideal^2)\)이다. 따라서
\[
        \mpideal^{\sigma-1}=(\mo c^{\sigma-1},\mpideal^\sigma)
\]
를 가정하고 양변에 \(\mpideal\)을 곱하면 분배법칙에 따라
\[
\begin{aligned}
        \mpideal^\sigma
        &=(\mpideal c^{\sigma-1},\mpideal^{\sigma+1})\\
        &=(\mo c^\sigma,\mpideal^2c^{\sigma-1},\mpideal^{\sigma+1})\\
        &=(\mo c^\sigma,\mpideal^{\sigma+1})
\end{aligned}
\]
이다.

증명할 보조정리는 다음 두 번째 형태로 바꿀 수 있다.
\[
\begin{gathered}
        \mq\equiv0(\mpideal^\sigma),\qquad
        \mq\not\equiv0(\mpideal^{\sigma+1}),\\
        \mpideal^{\sigma+\lambda}\equiv0(\mq),
        \qquad\lambda\ge1
\end{gathered}
\]
이면
\[
        \mpideal^{\sigma+\lambda-1}\equiv0(\mq)
\]
이기도 하다. 실제로 \(\mq\equiv0(\mpideal)\)이고 \(\mq\not\equiv0(\mpideal^{\sigma+1})\)이므로 모든 으뜸아이디얼에 그러한 지수 \(\sigma\ge1\)가 존재한다. 여기서 \(\mpideal^\varrho\equiv0(\mq)\)이고 1의 보조정리에 따라 \(\mpideal^\varrho\ne\mpideal^{\varrho+1}\)이므로 \(\mq\not\equiv0(\mpideal^{\varrho+1})\)이다. 두 번째 형태를 유한번 적용하면 \(\mpideal^\sigma\equiv0(\mq)\), 따라서 \(\mq=\mpideal^\sigma\)가 된다.

두 번째 형태의 가정과 \(\mpideal^\sigma\)의 위 표현을 함께 쓰면, 모든 \(q\in\mq\)에 대해
\[
        q\equiv bc^\sigma\pmod{\mpideal^{\sigma+1}},
        \qquad b\not\equiv0(\mpideal),
\]
또는
\[
        bc^\sigma\equiv0(\mq,\mpideal^{\sigma+1})
\]
이다. \((\mq,\mpideal^{\sigma+1})\)가 으뜸이므로 여기서
\[
        c^\sigma\equiv0(\mq,\mpideal^{\sigma+1})
\]
또는
\[
        c^{\sigma+\lambda-1}\equiv0(\mq,\mpideal^{\sigma+\lambda})
\]
가 따른다. 따라서 \(\mpideal^{\sigma+\lambda-1}\equiv0(\mq)\)이다.

4. 요약하면 다음과 같다.

\textbf{정리 VI.} \emph{가환환에서 공리 I--V가 성립하면, 영아이디얼 및 단위아이디얼과 다른 모든 아이디얼은 영아이디얼 및 단위아이디얼과 다른 유한개의 소아이디얼의 거듭제곱곱으로 유일하게 표현된다. 이 소아이디얼들은 단위아이디얼과 다른 진약수를 갖지 않는다.}

\begin{center}
{\bfseries\large § 9.\par}
\vspace{0.35em}
{\bfseries\large 가정한 분해에서 공리들을 이끌어 내기.\par}
\end{center}

정리 VI에 따라 공리 I--V에서 나오는 통상적인 아이디얼 분해의 존재로부터 역으로 이 공리들을 이끌어 내기 위해서는 분해를 다음 형태로 가정해야 한다.

\textbf{가정.} \emph{가환환 \(\mR\)에서 영원소 또는 단위원에서 생성되지 않는 모든 아이디얼은 소아이디얼들의 거듭제곱곱으로 유일하게 표현된다. 이 소아이디얼들은 단위원에서 생성되지 않는 단순아이디얼, 곧 \(\mo\)와 다른 진약수를 갖지 않는 아이디얼이다. 역으로 모든 단순아이디얼은 소아이디얼이다. 유일성은 영원소 또는 단위원에서 생성되지 않는 모든 아이디얼에 대해}
\[
        \ma^\varrho\ne\ma^{\varrho+1}
\]
\emph{이라는 강한 형태로 성립한다.}

여기서는 단위원의 존재를 가정하지 않는다. 단위원이 없으면 \glqq{}단위원에서 생성되지 않는다\grqq{}는 조건은 아무 제한도 주지 않는다.

\textbf{1. 공리 III(단위원의 존재)의 증명.} 공리 III이 성립하지 않는다고 하자. 그러면 가정에 어긋나게 \(\mo^2\)이 소아이디얼은 아니지만 단순아이디얼임을 보이면 된다. 강한 유일성 가정에 따라 \(\mo\ne\mo^2\)이다. 따라서 \(\mo\not\equiv0(\mo^2)\)이지만 \(\mo^2\equiv0(\mo^2)\)이므로 \(\mo^2\)은 소아이디얼이 아니다. 그러나 \(\mo^2\)은 단순하다. \(\mo^2\)은 \(\mo\)와 다른 어느 소아이디얼로도 나누어질 수 없으므로, 가정한 곱표현에 따라 \(\mo,\mo^2\)과 다른 약수를 갖지 않기 때문이다.

\textbf{2. 공리 IV(영인자가 없는 환)의 증명.} \(a\)가 영인자, 곧 \(ab=0\)이지만 \(a\ne0\), \(b\ne0\)이라고 하자. \(a,b\)에서 생성되는 주아이디얼의 표현들을 곱하면
\[
        (0)=\mpideal_1^{\varrho_1}\cdots\mpideal_k^{\varrho_k}
\]
를 얻는다. 단위원이 존재하므로 \(\mpideal_i\)들은 영아이디얼 및 단위아이디얼과 다르다. 어떤 한 인자를 생략하면 영아이디얼의 진약수가 되도록 이 표현이 최단이라고 가정한다. 영아이디얼에는 유일표현에 관한 가정을 두지 않았으므로 이 조건은 저절로 성립하는 것이 아니다. 이 표현에 소아이디얼이 하나만 나타나면
\[
        (0)=\mpideal^\varrho=\mpideal^{\varrho+1}
\]
이어서 강한 유일성 조건에 모순된다. 둘 이상의 소아이디얼이 나타나면
\[
\begin{aligned}
\mpideal_1^{\varrho_1+1}
 &=\bigl(\mpideal_1^{\varrho_1+1},
          \mpideal_1^{\varrho_1}\cdots\mpideal_k^{\varrho_k}\bigr)\\
 &=\mpideal_1^{\varrho_1}
   \bigl(\mpideal_1,\mpideal_2^{\varrho_2}\cdots\mpideal_k^{\varrho_k}\bigr)\\
 &=\mpideal_1^{\varrho_1}.
\end{aligned}
\]
단위원이 존재하므로 \(\mpideal_1\)은 나머지 모든 소아이디얼과 서로소이기 때문이다. 이 경우에도 강한 유일성 조건에 모순된다.

\textbf{3. 공리 I, II(사슬정리)의 증명.} 가정들로부터 모든 비영아이디얼에 대해 \emph{나눗셈에서 곱표현이 따른다}. 즉 \(\mm\equiv0(\ma)\), \(\ma\ne\mo\)이면 \(\mm=\ma\mb\)이면서 \(\mb\not\equiv0(\mm)\)인 \(\mb\)가 존재한다. 실제로
\[
        \mm\equiv0(\ma),\qquad
        \mm=\mpideal_1^{\varrho_1}\cdots\mpideal_r^{\varrho_r},\qquad
        \ma=\overline{\mpideal}_1^{\sigma_1}\cdots
             \overline{\mpideal}_{\bar r}^{\sigma_{\bar r}}
\]
라 하자. 그러면 각 \(\overline{\mpideal}\)은 어떤 \(\mpideal\)과 같고 \(\sigma_i\le\varrho_i\)이다. 따라서 일부 \(\sigma_i\)는 0일 수 있게 하여
\[
        \mb=\mpideal_1^{\varrho_1-\sigma_1}\cdots
             \mpideal_r^{\varrho_r-\sigma_r}
\]
로 놓으면 주장에 이른다. 그러므로 \(\mm\)의 약수들은 \(0\le\sigma_i\le\varrho_i\)의 조합들에 대응하는 유한개뿐이다. 따라서 \(\mm\)에 따른 잉여류환에는 이중사슬정리가 성립한다. 이로써 공리 I, II가 증명되었다. 약수사슬정리는 \(\mR\) 자체에도 성립한다. 영아이디얼의 모든 진약수는 영아이디얼과 다르기 때문이다.

\emph{분수체 안에서 정수적으로 닫혀 있다는 공리 V를 증명하려면 아이디얼 분해에서 나오는 통상적인 결론들이 필요하므로, 이를 먼저 유도해야 한다.}

4. \emph{모든 비영아이디얼에 따른 잉여류환은 주아이디얼환이다.} 단위아이디얼에 따른 잉여류환은 영원소만으로 이루어져 조건이 분명 성립하므로, 여기서 아이디얼은 단위아이디얼과도 다르다고 가정해도 된다.

먼저 아이디얼이 으뜸인 경우 \(\mq=\mpideal^\varrho\)라 하자. \(c\equiv0(\mpideal)\), \(c\not\equiv0(\mpideal^2)\)이면 3에 따라 \(\mo c=\mpideal\ma\), \((\ma,\mpideal)=\mo\)인 \(\ma\)가 존재한다. 따라서 모든 \(\sigma\le\varrho\)에 대해
\[
        \mpideal^\sigma=(\mo c^\sigma,\mpideal^\varrho)
\]
이고, 으뜸아이디얼에 따른 잉여류환의 모든 아이디얼은 주아이디얼이다. 일반적으로
\[
        \mm=\mq_1\mq_2\cdots\mq_r,
        \qquad \mq_i=\mpideal_i^{\varrho_i},
\]
이고 잉여류환이 이에 따라 직접합
\[
        \mR\mid\mm=\overline{\ma}_1+\cdots+\overline{\ma}_r
\]
으로 표현된다고 하자. \(\overline{\ma}_i\)는 \(\mR\mid\mq_i\)와 동형이므로(§ 4, 5) \(\overline{\ma}_i\)의 모든 아이디얼은 주아이디얼이다. \(\overline{\mc}\)가 잉여류환의 임의의 아이디얼이면 \(\overline{\ma}_i\overline{\mc}\)는 주아이디얼 \(\overline{\ma}_i\overline c_i\)이다. 따라서
\[
\begin{aligned}
        \overline{\mc}
        &=\overline{\mo}\,\overline{\mc}\\
        &=\overline{\ma}_1\overline{\mc}+\cdots+
          \overline{\ma}_r\overline{\mc}\\
        &=\overline{\ma}_1\overline c_1+\cdots+
          \overline{\ma}_r\overline c_r\\
        &=\overline{\ma}_1\overline c+\cdots+
          \overline{\ma}_r\overline c\\
        &=\overline{\mo}\,\overline c,
\end{aligned}
\]
여기서 \(\overline c=\overline c_1+\cdots+\overline c_r\)로 놓았다. 실제로 \(\overline{\ma}_i\overline{\ma}_k=(0)\)이고 \(\overline c_i\equiv0(\overline{\ma}_i)\)이므로, \(\overline{\ma}_i\) 안의 주아이디얼 \(\overline{\ma}_i\overline c_i\)와 \(\overline{\ma}_i\overline c\)는 일치한다.

잘 알려져 있듯이 주아이디얼환에 관한 정리는 다음 형태로도 쓸 수 있다.

\emph{\(\mc\)가 임의의 아이디얼이면, 주어진 한 아이디얼 \(\mb\)와 서로소인 아이디얼을 \(\mc\)에 곱하여 주아이디얼로 만들 수 있다.}

실제로 \(\mm=\mb\mc\)로 놓으면 \(\mc\)는 \(\mm\)을 법으로 주아이디얼이 된다. 즉
\[
        \mc=(\mo c,\mm)=(\mc\ma,\mc\mb)=\mc(\ma,\mb).
\]
따라서 \(\mo c=\mc\ma\)이고 \((\ma,\mb)=\mo\)이다.

\textbf{5. 분수아이디얼의 이론.} 분수체 \(\mK\) 안의 모든 유한 \(\mR\)-모듈을 분수아이디얼이라 한다.

\textbf{정리.} \emph{\(\mK\) 안의 비영 유한 \(\mR\)-모듈들은 곱셈에 관하여 아벨군을 이룬다.}

두 유한 \(\mR\)-모듈의 곱은 다시 유한 \(\mR\)-모듈이고, 곱셈은 결합적이며 가환적이다. \(\mA=\mo\mA\)이므로 단위아이디얼은 이 계의 단위원이다. 따라서 방정식 \(\mathfrak A\mX=\mo\)가 유한 \(\mR\)-모듈의 계 안에서 언제나 해를 갖는다는 것만 보이면 된다.

\textbf{예비주석.} 방정식 \(\mathfrak A\mT=\mathfrak B\)가 정확히 하나의 해를 가지면 \(\mT\)는 모듈 몫 \(\mathfrak B:\mathfrak A\)와 같다(§ 5, 4). 실제로
\[
\mT\equiv0(\mathfrak B:\mathfrak A),\qquad
\mathfrak B=\mathfrak A\mT
\equiv0\bigl(\mathfrak A\cdot(\mathfrak B:\mathfrak A)\bigr)
\equiv0(\mathfrak B),
\]
따라서
\[
        \mathfrak A\cdot(\mathfrak B:\mathfrak A)=\mathfrak B
\]
이다.

먼저 \(\mathfrak A\)가 주모듈 \(\mathfrak A=(\alpha)\)이면 \(\mX=(e/\alpha)\)가 \(\mathfrak A\mX=\mo\)의 해이고, \(\mX\)는 다시 유한 \(\mR\)-모듈이다. 이로부터 다음이 따른다. \emph{모든 유한 \(\mR\)-모듈은 \(\mR\)의 두 아이디얼의 모듈 몫이다.} 이 사실은 분수아이디얼이라는 이름을 정당화한다. 실제로 \(\tau_1,\ldots,\tau_r\)가 \(\mT\)의 모듈기저이고 \(\tau_i=t_i/a\)이며 \(\mt\)가 \(t_i\)들에서 생성되는 \(\mR\)의 아이디얼이면
\[
        \mo a\cdot\mT=\mt.
\]
따라서 예비주석에 따라 \(\mR\) 안에서 몫을 취하여 \(\mT=(\mt:\mo a)\)이다.

이제 \(\mathfrak A\mX=\mo\)의 해를 일반적으로 줄 수 있다. \(\mathfrak A=\ma:\mo c\)로 놓고 \(\ma\mb\)가 주아이디얼 \(\mo a\)가 되도록 \(\mb\)를 택한다. 그러면
\[
        \mathfrak A c=\ma,
        \qquad
        \mathfrak A\mb c=\mo a,
        \qquad
        \mathfrak A\cdot(\mb c:\mo a)=\mo.
\]
여기서 \(\mX\)는 한 아이디얼을 한 주아이디얼로 나눈 몫이므로 다시 유한 \(\mR\)-모듈이다. 이로써 \emph{군의 성질이 증명}되었다. 또 이로부터 \emph{임의의 두 아이디얼의 모듈 몫 \(\ma:\mb\)}는 방정식 \(\mb\mX=\ma\)의 해로서 \emph{유한 \(\mR\)-모듈}이라는 것이 따른다.

6. 군의 성질에서 다음 결론들이 나온다.

6\(\alpha\). \emph{약분할 수 있다.} 즉 \(\mathfrak T=\mathfrak A\mathfrak M:\mathfrak B\mathfrak M\)이면 \(\mathfrak T=\mathfrak A:\mathfrak B\)이고 그 역도 성립한다. 실제로 \(\mathfrak T\mathfrak B\mathfrak M=\mathfrak A\mathfrak M\)에서 \(\mathfrak T\mathfrak B=\mathfrak A\)가 나오며 그 역도 성립한다.

6\(\beta\). \emph{모든 분수아이디얼, 곧 유한 \(\mR\)-모듈은 \(\ma,\mb\)가 서로소인 표현 \(\mathfrak C=\ma:\mb\)를 갖고, 이 조건으로 \(\ma,\mb\)는 유일하게 결정된다.} 표현의 가능성은 5와 6\(\alpha\)에서 나온다. 한편
\[
        \mathfrak C=\ma:\mb=\overline{\ma}:\overline{\mb}
\]
이면 \(\ma\overline{\mb}=\mb\overline{\ma}\)이다. \(\ma,\mb\)도 서로소이고 \(\overline{\ma},\overline{\mb}\)도 서로소라고 가정하면 여기서 유일성이 나온다.

6\(\gamma\). \emph{정수적이지 않은 원소, 곧 \(\mK\)에는 속하지만 \(\mR\)에는 속하지 않는 원소에서 생성되는 모든 주모듈 \(\mo\eta\)는 두 주아이디얼의 몫으로 표현되되 분자의 어떤 거듭제곱도 분모로 나누어지지 않는다.} 기약표현 \(\mo\eta=\mb:\mc\), \((\mb,\mc)=\mo\)를 택하고, 4에 따라 \(\mo b=\mb\ma\), \((\ma,\mc)=\mo\)인 \(\ma\)를 택하자. 그러면
\[
        \mo\eta=\ma\mb:\ma\mc=\mo b:\ma\mc.
\]
또 \(\ma\mc=\mo b:\mo\eta\)이므로 \(\ma\mc\)도 주아이디얼이고, 따라서
\[
        \mo\eta=\mo b:\mo c.
\]
\((\ma,\mc)=\mo\), \((\mb,\mc)=\mo\)이므로 \(\ma\mb\)의 어떤 거듭제곱도 \(\mc\)로 나누어지지 않는다.

7. \emph{분수체 안에서 정수적으로 닫혀 있다는 공리 V의 증명.} 6\(\gamma\)에서 \emph{\(\mK\)의 모든 정수적이지 않은 원소는}
\[
        \eta=a/c
\]
\emph{로 표현되되, \(\mR\)에서 분자의 어떤 거듭제곱도 분모로 나누어지지 않게 할 수 있다.} 실제로 \(\mo\eta=\mo b:\mo c\)이면 \(\eta\)는 몫 \(b/c\)와 \(\mR\)의 한 단위원만큼만 다르다. 따라서 그러한 \(\eta\)는 \(\mR\)에 대하여 정수적임을 특징짓는 어떤 방정식도 만족할 수 없다. 실제로
\[
        \eta^n+r_1\eta^{n-1}+\cdots+r_n=0
\]
이면 \(\mR\)에서 \(a^n\equiv0(\mo c)\)가 되기 때문이다.

8. \emph{따름정리: 가환환 \(\mR\)에서 공리 I--IV가 성립하고 모든 으뜸아이디얼이 기약이면, 분수체 안에서 정수적으로 닫혀 있다는 공리 V도 성립한다.} 공리 I--III에 따라 모든 소아이디얼 거듭제곱 \(\mpideal^\varrho\)는 으뜸아이디얼이다. 특히 \(\mpideal^2\)도 그러하고, 가정에 따라 기약이다. 이제
\[
        \mpideal=(\mo c,\mpideal^2)
\]
임을 보이면 된다. 그러면 § 8, 3의 보조정리에 따라 모든 으뜸아이디얼이 소아이디얼 거듭제곱이 되고, 이는 다른 가정들과 함께 7에 따라 정수적 닫힘을 낳는다.

약수사슬정리에 따라
\[
        \mpideal=(\mo c_1,\ldots,\mo c_k,\mpideal^2)
\]
이고, 여기서 \(c_i\)의 곱셈자로는 \(\mpideal\)에 따른 잉여류들만 고려하면 되며, \(c_i\)들은 \(\mpideal\)에 따른 잉여류체에 관하여 선형독립이라고 가정할 수 있다. 이 선형독립성으로부터 \(k>1\)이면
\[
        \mpideal^2=[\mq_1,\mq_2],\qquad
        \mq_1=(\mo c_1,\mpideal^2),\qquad
        \mq_2=(\mo c_2,\ldots,\mo c_k,\mpideal^2)
\]
인 표현이 나온다. 이때 \(\mq_1,\mq_2\)는 \(\mpideal^2\)의 진약수이므로 \(\mpideal^2\)이 가약이 되어 가정에 모순된다. 이로써 \(\mpideal=(\mo c,\mpideal^2)\)와 그에 따르는 결론이 증명되었다.

역으로 공리 I--V의 결과로 모든 으뜸아이디얼이 기약이라는 것은 곧바로 알 수 있다. 소아이디얼 거듭제곱 \(\mpideal^\varrho\)는 \(\mpideal^\sigma,\mpideal^\tau\) 꼴의 진약수들의 최소공배수로 표현될 수 없기 때문이다.

9. \emph{환에 영인자를 허용하면, 공리 I--III과 몫환 안에서의 정수적 닫힘으로부터 모든 으뜸아이디얼이 기약이라는 것은 따르지 않는다}\footnote{각주 19를 보라.}. 실제로 \(\mT\)를 정수적 닫힘을 제외한 공리 I--IV가 모두 성립하는 환이라 하자. 그런 환은 존재한다. 예컨대 \(\mR\)에 공리 I--V가 성립할 때, \(\mR\)의 분수체의 한 유한 확대체 안에서 \(\mR\)에 대하여 정수적인 모든 원소의 계와 다른 유한 오더들이 그러하다(§ 3, 2). 따름정리 8에 따라 \(\mT\)에는 가약 으뜸아이디얼 \(\mq\)가 존재한다. \(\mpideal^\varrho\)를 \(\mq\)의 진배수라 하고 \(\overline{\mR}\)을 잉여류환 \(\mT\mid\mpideal^\varrho\)라 하자. 그러면 \(\mq\)에 대응하는 아이디얼 \(\overline{\mq}\)는 \(\overline{\mR}\)에서 비영 가약 으뜸아이디얼이 된다. \(\overline{\mR}\)에는 공리 I--III이 성립하고 몫환 안에서의 정수적 닫힘도 성립한다. 실제로 \(\mT\)에서 \(\mpideal\)로 나누어지지 않는 모든 원소는 \(\mpideal^\varrho\)와 서로소이므로, \(\overline{\mR}\)의 모든 정칙원소는 단위원이다. 따라서 \(\overline{\mR}\)은 그 몫환과 같다.

\begin{center}
{\bfseries\large § 10.\par}
\vspace{0.35em}
{\bfseries\large 이중사슬정리와 합성열.\par}
\end{center}

이하에서는 정렬되어 있다고 가정한 임의의 모듈영역(§ 2)에 대하여 \emph{이중사슬정리가 성립한다는 가정}, 곧 모든 모듈 약수사슬과 배수사슬이 유한 단계에서 멈춘다는 가정이 \emph{합성열이 존재한다는 가정}과 같음을 보인다\footnote{이 영역의 모듈들은 덧셈에 관하여 아벨군을 이룬다. 곱셈자영역이 \(\mR\)의 원소들로 이루어지므로 정확히는 \emph{일반화된} 아벨군이다. 아벨군이므로 합성열과 주열이라는 개념은 일치한다. 이 절의 정리들은 \glqq{}모듈\grqq{}을 임의의 비가환 일반화군을 포함한 한 군의 모든 정규부분군의 전체로 이해해도 성립한다. 뒤의 표기가 앞의 표기와 조금 다른 것은 군론을 떠올리게 하려는 것이다.}. 모듈영역을 가환환으로 특수화하면 환의 모든 아이디얼계에 대한 대응 사실들이 나온다. 이때 2에서는 줄곧 모듈동형을 환동형으로 바꾸어야 한다.

모듈영역 \(\mG\) 안에 \(\mG\) 자체와 영모듈 \(\mE\) 이외의 \(\mR\)-모듈이 없으면 \(\mG\)를 \emph{단순}하다고 한다. \(\mG\mid\mA\)가 단순하면 모듈 \(\mA\)를 \(\mG\) 안에서 \emph{단순}하다고 한다(극대 정규부분군).

\emph{배수사슬}
\[
        \mG,\mA_1,\ldots,\mA_r,\mE
\]
\emph{의 모든 모듈이 서로 다르고 각 모듈이 앞선 모듈 안에서 단순하면 이를 길이 \(r\)인 \(\mG\)의 합성열이라 한다.} 즉 잉여류모듈, 곧 몫군 \(\mA_{i-1}\mid\mA_i\)들과 \(\mG\mid\mA_1\), \(\mA_r\)가 단순 모듈영역이다. 잉여류모듈 \(\mG\mid\mA\)를 만들면 \(\mG\)에서 \(\mA\ne\mE\)까지의 합성열이 존재하는 경우를 합성열 자체의 경우로 환원할 수 있다.

1. \emph{\(\mG\)에서 이중사슬정리가 성립한다고 가정하면 \(\mG\)에는 합성열이 존재한다.} 여기서 \(\mG\ne\mE\)라고 가정한다. \(\mG\)의 정렬순서와 약수사슬정리로부터 § 6에서처럼 준사전식 순서를 써서 모든 모듈계의 정렬순서를 얻는다. 이 정렬순서와 약수사슬정리의 가정에 따라 \(\mG\) 안에서 단순한 모듈이 적어도 하나 존재한다. 실제로 \(\mG_1\)을 정렬순서에서 \(\mG\)의 첫 진약수라 하고, 일반적으로 \(\mG_i\)를 \(\mG_{i-1}\)의 첫 진약수라 하면 잘 정의된 사슬
\[
        \mG,\mG_1,\ldots,\mG_i,\ldots
\]
은 유한 단계에서 멈춘다. 따라서 반드시 \(\mG\) 안에서 단순한 모듈 \(\mA_1\)에 이른다. \(\mA_1\ne\mE\)이면 같은 방법으로 \(\mA_1\) 안에서 단순한 모듈 \(\mA_2\)를 얻는다. 일반적으로 \(\mA_{i-1}\ne\mE\)이면 \(\mA_{i-1}\) 안에서 단순한 모듈 \(\mA_i\)가 존재한다. 이 방식으로 잘 정의된 배수사슬
\[
        \mG,\mA_1,\ldots,\mA_i,\ldots
\]
을 얻는다. 이 사슬은 배수사슬정리의 가정에 따라 유한 단계에서 멈추므로 \(\mG\)의 합성열을 준다.

2. \emph{\(\mG\)에 합성열이 존재한다고 가정하면 \(\mG\)에서 이중사슬정리가 성립한다.} 증명은 완전귀납법으로 하며, \(\mG\)에 정렬순서를 가정할 필요는 없다. 귀납과정에서는 합성열의 존재만을 가정한 조르당--횔더 정리의 증명도 함께 나온다\footnote{환의 경우는 Masazo Sono(각주 19), \emph{On Congruences I}, §§ 11--13을 보라. 그곳에서는 비가환군에서 합성열에 대응하는 것, 곧 \(\mR,\mA_1,\ldots,\mA_r,\mE\) 꼴의 열도 다룬다. 여기서 각 \(\mA_i\)는 \(\mA_{i-1}\) 안의 단순아이디얼이지만 \(\mR\)의 아이디얼일 필요는 없다. 실제로 비가환군에서도 배수사슬을 각각 \(\mA_i\)가 \(\mA_{i-1}\)의 정규부분군인 \(\mG,\mA_1,\ldots,\mA_r,\ldots\)로 이해하고 약수사슬도 그에 맞게 정의하면 위 논의가 모두 유지된다. 또 Dedekind, \emph{Über die von drei Moduln erzeugte Dualgruppe}, Math. Ann. 53 (1900), 371--403쪽을 보라. 그 논문에서도 훨씬 일반적인 영역(모듈군)에 대해 존재만을 가정하여 조르당--횔더 정리를 증명한다(§ 6, XVI). 여기서는 제2동형정리 대신 조금 더 약한 대응관계가 나타나므로 정리의 진술도 조금 더 약하지만, 증명은 원리상 위의 증명과 같다.}.

먼저 \(\mG\)가 단순, 곧 \(r=0\)이어서 합성열 \(\mG,\mE\)만 존재한다고 하자. 이때 다음 명제들이 성립한다.
\[
\begin{array}{ll}
\alpha) & \text{\(\mG\) 안에서 단순한 모든 모듈을 지나는 합성열을 그을 수 있다.}\\
\beta) & \text{조르당--횔더 정리: 모든 합성열의 길이는 같고, 몫군}\\
       & \text{(잉여류모듈)의 계는 순서를 제외하고 일치한다. 즉 대응하는}\\
       & \text{몫군들은 동형이다.}\\
\gamma) & \text{\(\mG\)와 다른 임의의 모듈을 지나는 합성열을 그을 수 있다.}\\
\delta) & \text{배수사슬정리가 성립한다.}\\
\varepsilon) & \text{약수사슬정리가 성립한다.}
\end{array}
\]
따라서 길이 \(\bar r<r+1\)인 합성열을 갖는 모든 모듈영역에는 \(\alpha)\)--\(\varepsilon)\)이 성립한다고 가정할 수 있다. 이제 \emph{\(\mG\)에 길이 \(r+1\)인 합성열}
\[
        \mG,\mA,\mA_1,\ldots,\mA_r,\mE
\]
\emph{이 존재한다는 가정} 아래에서 이 명제들을 증명한다.

\textbf{2\(\alpha\).} \(\mA\)와 다른, \(\mG\) 안에서 단순한 모듈이 없으면 증명할 것이 없다. \(\mB\ne\mA\)가 \(\mG\) 안에서 단순하다고 하자. \(\mB\)를 지나는 \(\mG\)의 합성열이 존재함을 보여야 한다. \(\mA,\mB\)가 모두 \(\mG\) 안에서 단순하고 서로 다르므로 반드시 \((\mA,\mB)=\mG\)이다. \(\mM=[\mA,\mB]\)로 놓으면 제2동형정리(§ 4, 2)에 따라
\[
        \mG\mid\mB\simeq\mA\mid\mM,
        \qquad
        \mG\mid\mA\simeq\mB\mid\mM.
\]
따라서 \(\mM\)은 \(\mA,\mB\) 안에서 단순하다. \(\mA\)는 길이 \(r<r+1\)인 합성열을 가지므로 귀납가정 \(\alpha)\), \(\beta)\)에 따라 \(\mA\)에는 \(\mM\)을 지나는 길이 \(r\)인 합성열, 이를테면
\[
        \mA,\mM,\mM_1,\ldots,\mM_{r-1},\mE
\]
이 존재한다. 그러므로
\[
        \mG,\mB,\mM,\mM_1,\ldots,\mM_{r-1},\mE
\]
는 \(\mB\)를 지나는 \(\mG\)의 합성열이다.

\textbf{2\(\beta\).} 조르당--횔더 정리를 증명하기 위하여
\[
\begin{aligned}
(1)\quad&\mG,\mA,\mA_1,\ldots,\mA_r,\mE,\\
(2)\quad&\mG,\mB,\mB_1,\ldots,\mB_s,\mE,
\qquad s\ge r
\end{aligned}
\]
를 \(\mG\)의 두 합성열이라 하자. \(\mA=\mB\)이면 \(r<r+1\)에 대한 귀납가정으로 모두 끝난다. 그렇지 않으면 2\(\alpha\)에서 구성한 열
\[
\begin{aligned}
(3)\quad&\mG,\mA,\mM,\mM_1,\ldots,\mM_{r-1},\mE,\\
(4)\quad&\mG,\mB,\mM,\mM_1,\ldots,\mM_{r-1},\mE
\end{aligned}
\]
과 비교한다. \(r<r+1\)에 대한 귀납가정에 따라 (1)과 (3)의 몫군계가 일치하고, (4)에서 \(\mB\)를 지나는 합성열의 길이가 \(r\)이므로 (2)와 (4)의 몫군계도 일치한다. 따라서 \(s=r\)이다. 또한 2\(\alpha\)의 동형에 따라 (3)과 (4)가 일치한다. 이로써 (1)과 (2)가 일치하며 조르당--횔더 정리가 증명되었다.

\textbf{2\(\gamma\).} 예비주석. \emph{\(\mA,\mB,\mH\)가 \(\mG\)의 모듈이고 \(\mB\)가 \(\mA\) 안에서 단순하면, \((\mB,\mH)\)와 \((\mA,\mH)\)는 서로 같거나 \((\mB,\mH)\)가 \((\mA,\mH)\) 안에서 단순하다.} 제2동형정리(§ 4, 2)에 따라
\[
        (\mA,\mB,\mH)\mid(\mB,\mH)
        \simeq \mA\mid[\mA,(\mB,\mH)]
\]
이다. \(\mB\equiv0(\mA)\)이므로 이는
\[
        (\mA,\mH)\mid(\mB,\mH)\simeq\mA\mid\mC,
        \qquad
        \mC=[\mA,(\mB,\mH)]
\]
로 쓸 수 있다. 여기서 \(\mC\equiv0(\mA)\)이다. 한편 \(\mB\equiv0(\mA)\), \(\mB\equiv0((\mB,\mH))\)이므로 \(\mB\equiv0(\mC)\)이기도 하다. 따라서 \(\mC\)는 \(\mA,\mB\) 사이에 놓이고, 가정에 따라 \(\mA\) 또는 \(\mB\)와 같아 예비주석이 따른다.

이제 다시 \(\mG,\mA,\mA_1,\ldots,\mA_r,\mE\)를 \(\mG\)의 한 합성열이라 하고, \(\mH\ne\mG\)를 \(\mG\)의 임의의 모듈이라 하자. \(\mH\)를 지나는 합성열이 존재함을 보이기 위하여 모듈의 열
\[
        \mG,(\mA,\mH),(\mA_1,\mH),\ldots,(\mA_r,\mH),\mH
\]
을 만든다. 예비주석에 따라 이들 가운데 서로 다른 모듈 \(\mG,\mC,\mC_1,\ldots,\mC_s,\mH\)는 \(\mG\)에서 \(\mH\)까지의 합성열을 이룬다. 따라서 \(\mC\)는 \(\mG\) 안에서 단순한 모듈이고, 특수한 경우 2\(\alpha\)에서는 \(\mH\)와 같다. 그러므로 2\(\alpha\)에 따라 \(\mC\)에는 길이 \(r<r+1\)인 합성열이 존재하고, 귀납가정에 따라 \(\mC\)에서 \(\mH\)를 지나는 합성열을 그을 수 있다. 여기에 \(\mG\)를 덧붙이면 \(\mG\)에서 \(\mH\)를 지나는 합성열을 얻는다. 이 논의를 되풀이하면 특히 그러한 합성열을 \(\mC\)에서 \(\mH\)까지 구성한 열의 연장으로 택할 수 있음도 알 수 있다.

\textbf{2\(\delta\).} 배수사슬정리의 증명. 다시 \(\mG\)에 길이 \(r+1\)인 합성열이 존재한다고 가정하고,
\[
        \mG,\mB,\mB_1,\ldots,\mB_\nu,\ldots
\]
를 각 모듈이 앞선 모듈의 진배수인 배수사슬이라 하자. 언제나 \(\mG\)를 첫 항으로 덧붙일 수 있으므로 사슬이 \(\mG\)에서 시작한다고 해도 일반성을 잃지 않는다. 2\(\gamma\)에 따라 \(\mG\) 안에는 \(\mB\)를 지나는 합성열이 있고, 따라서 \(\mB\)는 더 짧은 합성열을 갖는다. 그러므로 귀납가정에 따라 사슬 \(\mB,\mB_1,\ldots,\mB_\nu,\ldots\)은 유한 단계에서 멈추고, \(\mG\)를 덧붙인 사슬도 마찬가지이다.

\textbf{2\(\varepsilon\).} 약수사슬정리의 증명. 다시 \(\mG\)에 길이 \(r+1\)인 합성열이 존재한다고 가정하고,
\[
        \mE,\mC,\mC_1,\ldots,\mC_\nu,\ldots
\]
를 각 모듈이 앞선 모듈의 진약수인 약수사슬이라 하자. 사슬이 \(\mE\)에서 시작한다고 해도 역시 일반성을 잃지 않는다. 2\(\gamma\)에 따라 \(\mG\) 안에는 \(\mC\)를 지나는 합성열이 있고, 따라서 \(\mG\mid\mC\)에는 더 짧은 합성열이 있다. 그러므로 귀납가정에 따라 \(\mG\mid\mC\)의 약수사슬
\[
        \mC\mid\mC,\mC_1\mid\mC,\ldots,\mC_\nu\mid\mC,\ldots
\]
은 유한 단계에서 멈춘다. 제1동형정리(§ 4, 2)에 따른 일대일 대응 때문에 원래의 \(\mC\)에서 시작하는 사슬도 마찬가지이고, \(\mE\)를 덧붙인 사슬도 그러하다.

\emph{이로써 두 가정, 곧 합성열의 존재와 이중사슬정리의 동치가 증명되었다.}

\begin{center}
(1925년 8월 13일 접수.)
\end{center}

\clearpage
\setcounter{footnote}{0}
